%*******************************************************************************
%****************************** Chapter 5 *************************************
%*******************************************************************************

\chapter{Agitation}
\label{ch:agitation}

\ifpdf
    \graphicspath{{Chapter5/Figs/Raster/}{Chapter5/Figs/PDF/}{Chapter5/Figs/}}
\else
    \graphicspath{{Chapter5/Figs/Vector/}{Chapter5/Figs/}}
\fi

This chapter covers the agitation application of MIXR-1, which serves the CHE 3214L experiment. It presents the speed-measurement chain, the motor drive and its closed-loop speed controller, the torque-sensing chain and the power computation derived from it, the two user-facing modes that operate the pathway, and the validation of each.


%********************************** 5.1 **************************************
\section{Agitation Subsystem Overview}
\label{sec:agit-overview}

% TODO: describe both impellers: type, diameter, number and width of blades, clearance from the vessel floor, the vessel diameter and
% liquid height, and the baffles. Chapter 3 and Chapter 4 do not mention the 174 mm impeller at all; add it where the geometry is given.
% TODO: Figure 5.1 -- agitation pathway block diagram: setpoint -> PI controller -> PWM ->
% driver -> motor -> shaft -> {encoder -> speed estimate -> feedback, torque sensor -> power},
% with the power limiter shown acting on the controller's output ceiling.

The apparatus is operated with two impellers, of 96~mm and of 174~mm diameter, which the firmware offers in a menu with different speed limits for the knob, 1400~rpm for the 96~mm impeller and 350~rpm for the 174~mm impeller. Both are driven by the same motor, sensor, and controller, so every quantity of this chapter that depends on the load, namely the plant of Section~\ref{subsec:agit-pi-tuning}, the closed-loop behaviour, and the torque and Power number, is determined for each of them.

The pathway closes one loop, on speed, and makes one open measurement, of torque. Both feed the $N_{Re}$--$N_{Po}$ computation: speed supplies $n$ in Equations~\ref{eq:reynolds} and~\ref{eq:power}, and torque supplies the shaft power $P$.


%********************************** 5.2 **************************************
\section{Rotary Encoder and Speed Measurement}
\label{sec:agit-speed}

This section covers the complete speed-measurement chain: the encoder and its configuration, the quadrature decoding performed in the interrupt path, the validation of the index channel, the speed estimator that converts timestamped counts into revolutions per minute, and the analysis of its error. The estimator follows the M/T family of methods reviewed in Section~\ref{subsec:rrl-encoder}, and the error analysis uses the notation of \citet{anuchin2019}.

\subsection{Sensor Selection and Interfacing}
\label{subsec:agit-encoder-sel}

The speed sensor is an AMT102-V modular capacitive incremental encoder~\citep{sameskyamt10}, mounted on the motor shaft. It outputs two quadrature channels, A and B, and a once-per-revolution index channel X, as CMOS square waves. Capacitive sensing removes the optical disc of conventional encoders, and the resolution is selected by a four-position DIP switch from sixteen values between 48 and 2048 pulses per revolution (PPR). The device characteristics relevant to this work are summarized in Table~\ref{tab:amt102}.

\begin{table}[htbp]
    \centering
    \caption[AMT102-V characteristics]{AMT102-V characteristics used in this work, from the manufacturer datasheet~\citep{sameskyamt10}.}
    \label{tab:amt102}
    \small
    \begin{tabular}{@{}p{0.26\textwidth}p{0.29\textwidth}p{0.37\textwidth}@{}}
        \toprule
        Parameter & Value & Condition \\
        \midrule
        Supply voltage & 3.6 -- 5.5 V (typ.\ 5 V) & $V_{DD}$ \\
        Supply current & 6 mA & unloaded outputs \\
        Output high level & $\geq V_{DD} - 0.8$ V & \\
        Output low level & $\leq 0.4$ V & \\
        Output current & 2 mA & per channel, sink or source \\
        Rise/fall time & 30 ns (typ.) & \\
        Channels & A, B (quadrature), X (index) & \\
        Phase difference & $90^{\circ}$ (typ.) & A leads B for CCW, viewed from front \\
        Index & one pulse per revolution, one state wide & \\
        Accuracy & $0.25^{\circ}$ & \\
        Duty cycle & 49 -- 51 \% & at 256, 512, 1024, 2048 PPR \\
        Maximum speed & 15\,000 rpm & at 256 PPR \\
        Operating temperature & $-40$ to $100\,^{\circ}$C & \\
        \bottomrule
    \end{tabular}
\end{table}

\paragraph{Resolution.}
The encoder is set to $K = 256$ PPR (DIP switches 1--4 set to 0, 0, 1, 1). The datasheet defines the counts per revolution as four times the PPR, which is the fourfold decoding of Equation~1 in \citet{anuchin2019}:
\begin{equation}
    N = 4K = 4 \times 256 = 1024 \text{ counts per revolution}
    \label{eq:cpr}
\end{equation}
Each count therefore corresponds to $360^{\circ}/N = 0.352^{\circ}$ of shaft rotation, and one full encoder period, which is the aligned increment used by the estimator in Section~\ref{subsec:agit-cet}, to $360^{\circ}/K = 1.406^{\circ}$. For comparison, the rated accuracy of the encoder is $0.25^{\circ}$, so the count resolution is of the same order as the encoder's own position accuracy; a finer setting would make the edges more frequent without making the position more accurate. The choice of 256 PPR among the sixteen available settings is analysed and tested in Section~\ref{subsec:agit-ppr}.

\paragraph{Electrical interface.}
The encoder outputs are connected to GPIO 24 (A), GPIO 23 (B), and GPIO 22 (X) of the Raspberry Pi through a four-channel bidirectional logic level converter (SparkFun, supplied by MakerLab, SKU MLE00152). The minimum supply of the encoder is 3.6~V, so it cannot be powered from the 3.3~V rail and is supplied at 5~V; at that supply its high output level is at least $V_{DD} - 0.8 = 4.2$~V, above the 3.3~V logic level of the Raspberry Pi inputs, which are not 5~V tolerant. The converter translates each signal between the two levels. Its high-voltage pin is connected to the 5~V encoder supply, its low-voltage pin to the 3.3~V rail of the Raspberry Pi, and its ground to the common ground. Channels A, B, and X use three of the four channels, leaving one spare. The firmware also enables the internal pull-up of each GPIO input, which is redundant with the pull-ups of the converter and harmless.

\paragraph{Effect of the converter on the signal timing.}
A converter of this type switches with a MOSFET and holds each side high through a pull-up resistor. Its falling edge is driven through the MOSFET and is fast, but its rising edge is set by the pull-up resistance $R$ and the load capacitance $C$, with a 10--90~\% rise time of about $2.2RC$. For a 10~k$\Omega$ pull-up and an assumed load of 50~pF for the GPIO input and the wiring, the rise time is about 1.1~$\mu$s. It has three consequences for the speed measurement.
\begin{enumerate}
    \item At 256 PPR the converter is not the limiting element. A rise time of about 1~$\mu$s is a fifth of the 5~$\mu$s sample interval of the pigpio sampler and under 5~\% of the 23.4~$\mu$s state width at 2500~rpm, so the sampler remains the limit on the edge rate, as analysed in Section~\ref{subsec:agit-ppr}.
    \item It adds an asymmetry between rising and falling edges. A rising edge delayed by about 1~$\mu$s against a falling edge shortens the apparent high time of each channel by that amount, which is about 1.2~\% of the 93.75~$\mu$s channel period at 256 PPR and 2500~rpm. This is of the same size as the $\pm 1$~\% duty-cycle tolerance of the encoder itself. The slot-aligned estimator cancels it, because each snapshot is stored at the same type of edge and a constant delay of one edge type cancels in the difference between two snapshots.
    \item The same absolute delay is a larger fraction of the period at a higher resolution. At 1024 PPR the channel period at 2500~rpm is 23.4~$\mu$s and the same 1.1~$\mu$s would be 4.7~\% of it, which adds a further reason, beyond the edge rate, not to select a higher resolution with this converter.
\end{enumerate}

% TODO: confirm the converter's pull-up value and transistor type from the SparkFun schematic of
% the board supplied (the 10 kOhm figure is the design value assumed here), and measure the rise
% and fall times and the duty cycle at the GPIO pins with an oscilloscope at one or two speeds.
% Replace the assumed 50 pF estimate with the measured rise time.
% TODO: state which rail supplies the 5 V of the encoder and of the converter's HV pin
% (Section~\ref{sec:plat-power}), and show the converter in the schematic
% (Appendix~\ref{app:schematics}). Add it to the BOM (Appendix~\ref{app:bom}), PHP 75.
% TODO: Figure 5.2 -- encoder mounting photograph showing the sleeve, shaft adapter and base.
% TODO: confirm positive count direction against a known rotation. The datasheet defines
% A-leads-B as counter-clockwise viewed from the front. In the preflight run of the Pi (preflight.csv) delta_count is negative
% in all 799 rows while raw_rpm is positive and revs counts up, which is only consistent with ENCODER_DIRECTION = -1 in the build
% that ran; the config.hpp supplied says +1. State the direction convention (which way the shaft turns when the speed is positive)
% and the value of ENCODER_DIRECTION that was used.

\subsection{Resolution Selection}
\label{subsec:agit-ppr}

The AMT102-V offers sixteen resolutions between 48 and 2048 PPR, selected by the DIP switch. The resolution is a design decision rather than a given, because it trades three quantities against one another: the speed resolution available at low speed, the edge rate that the Raspberry Pi must service at high speed, and the quality of the encoder signals, which the datasheet specifies differently for different settings. This section evaluates all sixteen settings against those criteria analytically, reports a measured comparison of eight of them, and states the basis for the choice of 256 PPR and the limits of that basis.

\paragraph{Selection criteria.}
Four criteria were applied.
\begin{enumerate}
    \item \emph{Decoder feasibility at top speed.} The decoder detects edges through the pigpio sampler, whose default sample interval is 5~$\mu$s, and level changes shorter than that interval may be missed~\citep{pigpio}. Requiring a state width of at least two sample intervals, the highest PPR that can be decoded at speed $n$ is
    \begin{equation}
        K_{\max}(n) = \frac{60}{4\, n\, S_{\min}}, \qquad S_{\min} = 10\ \mu\text{s}
        \label{eq:kmax}
    \end{equation}
    where $S_{\min}$ is the minimum state width. For the highest speed the apparatus reaches, taken as 2500~rpm, this gives $K_{\max} = 600$ PPR. Equivalently, for a given setting the highest decodable speed is $n_{\max} = \frac{60}{N S_{\min}}$, which for 256 PPR is 5860~rpm.
    \item \emph{Signal quality.} The datasheet specifies the duty cycle of the A and B channels as 49--51~\% at 256, 512, 1024, and 2048 PPR, 47--53~\% at 48, 96, 100, 125, 192, 200, 250, 384, 400, and 500 PPR, and 43--56~\% at 800 and 1000 PPR~\citep{sameskyamt10}. Duty-cycle error is one of the encoder non-idealities that corrupts an estimator not protected by slot alignment~\citep{anuchin2019}.
    \item \emph{Low-speed behaviour of the estimator.} The slot-aligned estimator of Section~\ref{subsec:agit-cet} produces a new estimate on every control-loop period only while the aligned period does not exceed the loop period, that is above $n_{\text{upd}} = \frac{60}{K\,T_s}$ (Equation~\ref{eq:nupd}), and it reports zero below $n_{\min} = \frac{60}{K \times 0.1\,\text{s}}$ (Equation~\ref{eq:nmin}). A higher PPR lowers both. For a lowest controlled speed $n_{\text{low}}$, refreshing the estimate every loop requires $K \geq \frac{60}{n_{\text{low}} T_s}$.
    \item \emph{Measured noise.} The standard deviation and coefficient of variation of the settled speed at each operating point, compared across settings.
\end{enumerate}

% TODO: confirm the 2500 rpm top speed from the new open-loop sweep (highest PWM step, no
% load) and update Equation~\ref{eq:kmax} and Tables~\ref{tab:ppr-spec} if it differs.

\paragraph{Analytical comparison of all sixteen settings.}
Table~\ref{tab:ppr-spec} gives the datasheet properties of every setting and its feasibility at 2500~rpm, and Table~\ref{tab:ppr-est} gives its consequences for the two estimators.

\begin{table}[htbp]
    \centering
    \caption[All AMT102-V resolution settings: datasheet and decoder feasibility]{All sixteen AMT102-V resolution settings, with the DIP switch code (switches 1--4), rated maximum speed, and duty-cycle tolerance from the datasheet~\citep{sameskyamt10}, and the edge rate and state width at 2500~rpm. A setting is decodable when the state width is at least two 5~$\mu$s sample intervals (Equation~\ref{eq:kmax}).}
    \label{tab:ppr-spec}
    \resizebox{\textwidth}{!}{%
    \begin{tabular}{@{}rrcrcrrrc@{}}
        \toprule
        PPR & CPR & DIP & Rated max (rpm) & Duty (\%) & Edge rate (kHz) & State width ($\mu$s) & Width / sample & Decodable \\
        \midrule
        48 & 192 & 1111 & 15{,}000 & 47--53 & 8.0 & 125.0 & 25.0 & yes \\
        96 & 384 & 1101 & 15{,}000 & 47--53 & 16.0 & 62.5 & 12.5 & yes \\
        100 & 400 & 0111 & 15{,}000 & 47--53 & 16.7 & 60.0 & 12.0 & yes \\
        125 & 500 & 1011 & 15{,}000 & 47--53 & 20.8 & 48.0 & 9.6 & yes \\
        192 & 768 & 1110 & 7{,}500 & 47--53 & 32.0 & 31.2 & 6.2 & yes \\
        200 & 800 & 0101 & 15{,}000 & 47--53 & 33.3 & 30.0 & 6.0 & yes \\
        250 & 1{,}000 & 1001 & 15{,}000 & 47--53 & 41.7 & 24.0 & 4.8 & yes \\
        256 & 1{,}024 & 0011 & 15{,}000 & 49--51 & 42.7 & 23.4 & 4.7 & yes \\
        384 & 1{,}536 & 1100 & 7{,}500 & 47--53 & 64.0 & 15.6 & 3.1 & yes \\
        400 & 1{,}600 & 0110 & 7{,}500 & 47--53 & 66.7 & 15.0 & 3.0 & yes \\
        500 & 2{,}000 & 1010 & 7{,}500 & 47--53 & 83.3 & 12.0 & 2.4 & yes \\
        512 & 2{,}048 & 0001 & 15{,}000 & 49--51 & 85.3 & 11.7 & 2.3 & yes \\
        800 & 3{,}200 & 0100 & 7{,}500 & 43--56 & 133.3 & 7.5 & 1.5 & no \\
        1000 & 4{,}000 & 1000 & 7{,}500 & 43--56 & 166.7 & 6.0 & 1.2 & no \\
        1024 & 4{,}096 & 0010 & 7{,}500 & 49--51 & 170.7 & 5.9 & 1.2 & no \\
        2048 & 8{,}192 & 0000 & 7{,}500 & 49--51 & 341.3 & 2.9 & 0.6 & no \\
        \bottomrule
    \end{tabular}}
\end{table}

\begin{table}[htbp]
    \centering
    \caption[All AMT102-V resolution settings: estimator consequences]{Consequences of each resolution setting for the two estimators. The frequency-based step is the speed represented by one count in a fixed window, $\frac{60}{T_s N}$. The update floor $n_{\text{upd}}$ and the minimum reportable speed $n_{\min}$ apply to the slot-aligned estimator with a 10~ms loop and a 100~ms stopped-shaft timeout.}
    \label{tab:ppr-est}
    \begin{tabular}{@{}rrrrr@{}}
        \toprule
        PPR & Step, 10 ms (rpm) & Step, 20 ms (rpm) & $n_{\text{upd}}$ (rpm) & $n_{\min}$ (rpm) \\
        \midrule
        48 & 31.25 & 15.62 & 125.0 & 12.50 \\
        96 & 15.62 & 7.81 & 62.5 & 6.25 \\
        100 & 15.00 & 7.50 & 60.0 & 6.00 \\
        125 & 12.00 & 6.00 & 48.0 & 4.80 \\
        192 & 7.81 & 3.91 & 31.2 & 3.12 \\
        200 & 7.50 & 3.75 & 30.0 & 3.00 \\
        250 & 6.00 & 3.00 & 24.0 & 2.40 \\
        256 & 5.86 & 2.93 & 23.4 & 2.34 \\
        384 & 3.91 & 1.95 & 15.6 & 1.56 \\
        400 & 3.75 & 1.88 & 15.0 & 1.50 \\
        500 & 3.00 & 1.50 & 12.0 & 1.20 \\
        512 & 2.93 & 1.46 & 11.7 & 1.17 \\
        800 & 1.88 & 0.94 & 7.5 & 0.75 \\
        1000 & 1.50 & 0.75 & 6.0 & 0.60 \\
        1024 & 1.46 & 0.73 & 5.9 & 0.59 \\
        2048 & 0.73 & 0.37 & 2.9 & 0.29 \\
        \bottomrule
    \end{tabular}
\end{table}

Three results follow from the tables. First, criterion~1 excludes 800, 1000, 1024, and 2048 PPR outright: at 2500~rpm their state width is 7.5, 6.0, 5.9, and 2.9~$\mu$s, between 0.6 and 1.5 sample intervals, so edges would be lost. The settings from 384 to 512 PPR pass but with narrowing margin, the width at 512 PPR being 11.7~$\mu$s. Second, the settings rated to 7500~rpm, which include 192 and 384 to 500 PPR, remain well above the 2500~rpm top speed, so the rated maximum does not discriminate among the decodable settings. Third, the low-speed quantities in Table~\ref{tab:ppr-est} improve in direct proportion to PPR: 48 PPR cannot refresh its estimate every loop below 125~rpm and reports zero below 12.5~rpm, whereas 256 PPR does so down to 23.4~rpm and reports speeds down to 2.34~rpm.

\paragraph{Measured comparison with the earlier estimator.}
The eight settings 48, 96, 100, 125, 192, 200, 250, and 256 PPR were compared on 2 September, at sampling windows of 10~ms and 20~ms, with the earlier discrete-time (DCT) estimator and before the slot-aligned estimator was adopted. The motor was driven open loop through PWM steps of 10~\% from 0 to 90~\%, each held for 10~s, with the speed logged at 100~Hz. Only the interior seconds of each step (1 to 9~s, 800 samples) were analysed, so that the first and last second of each step, which contain the command transient, are excluded. For each step the standard deviation of the raw speed was computed, and the spectral decomposition of Section~\ref{subsec:agit-fft} was used to separate it into the line at the rotation frequency and a broadband remainder. Table~\ref{tab:ppr-dct} and Figure~\ref{fig:ppr-noise} give the result.

\begin{table}[htbp]
    \centering
    \caption[Resolution comparison with the earlier estimator]{Resolution comparison with the earlier (DCT) estimator, from the open-loop PWM sweeps of 2 September. SD is the standard deviation of the raw speed, $A_1$ the amplitude of the line at the rotation frequency, and noise the root-mean-square value of the remainder (Section~\ref{subsec:agit-fft}). The first six value columns are means over the nine steps from 10 to 90~\% PWM; the last two are the noise at the 10~\% step, about 220~rpm. All values in rpm.}
    \label{tab:ppr-dct}
    \small
    \resizebox{\textwidth}{!}{%
    \begin{tabular}{@{}rrrrrrrrr@{}}
        \toprule
        & \multicolumn{2}{c}{SD} & \multicolumn{2}{c}{$A_1$} & \multicolumn{2}{c}{Noise} & \multicolumn{2}{c}{Noise at 220 rpm} \\
        \cmidrule(lr){2-3} \cmidrule(lr){4-5} \cmidrule(lr){6-7} \cmidrule(lr){8-9}
        PPR & 10 ms & 20 ms & 10 ms & 20 ms & 10 ms & 20 ms & 10 ms & 20 ms \\
        \midrule
        48 & 22.3 & 15.9 & 10.1 & 7.7 & 16.7 & 12.4 & 9.2 & 6.8 \\
        96 & 23.4 & 15.3 & 11.6 & 6.2 & 18.5 & 12.3 & 8.0 & 5.2 \\
        100 & 19.8 & 15.0 & 7.6 & 5.4 & 16.7 & 12.1 & 6.6 & 6.1 \\
        125 & 19.6 & 14.5 & 7.5 & 4.7 & 17.2 & 12.3 & 5.1 & 6.0 \\
        192 & 22.5 & 17.2 & 8.4 & 5.7 & 19.1 & 14.9 & 5.7 & 6.3 \\
        200 & 23.6 & 18.3 & 10.4 & 7.2 & 18.6 & 14.8 & 6.8 & 6.2 \\
        250 & 23.5 & 20.9 & 10.5 & 7.3 & 19.4 & 14.5 & 5.6 & 5.9 \\
        256 & 55.9 & 40.2 & 10.0 & 6.4 & 30.6 & 21.7 & 4.5 & 4.7 \\
        \bottomrule
    \end{tabular}}
\end{table}

\begin{figure}[htbp]
    \centering
    \includegraphics[width=\textwidth]{fig_ppr_noise}
    \caption[Noise and rotation-frequency line against PPR]{Left: broadband noise at the 10~\% step, about 220~rpm, against PPR for the two windows, with the counting-quantization prediction $\frac{60}{T_w N\sqrt{12}}$ for each window (dotted). Right: amplitude of the line at the rotation frequency, mean over the nine steps, against PPR.}
    \label{fig:ppr-noise}
\end{figure}

The measurements show the following.
\begin{itemize}
    \item \emph{The shaft-synchronous line does not depend on the resolution.} The amplitude of the line at the rotation frequency is between 4.7 and 11.6~rpm at every setting and shows no trend with PPR. It is a property of the shaft and not of the encoder resolution, which is the behaviour expected of an error that repeats every revolution (Section~\ref{subsec:agit-fft}).
    \item \emph{Counting quantization matters only at the lowest settings.} At about 220~rpm the broadband noise at 48 PPR and 10~ms is 9.2~rpm, which agrees with the quantization prediction of 9.0~rpm. It falls to 4.5 to 6.8~rpm for the settings from 100 PPR upward, above the prediction of 1.7 to 4.3~rpm for those settings, so a floor of about 4 to 6~rpm remains that resolution does not reduce.
    \item \emph{Averaged over the speed range, the broadband noise does not fall with PPR.} It is 12 to 19~rpm at every setting from 48 to 250 PPR. It is higher at 256 PPR, 22 to 31~rpm, because the runs at 250 and 256 PPR contain isolated raw samples more than twice the true speed at the highest steps, up to 6102~rpm and 7119~rpm at 10~ms against a top speed near 2500~rpm.
    \item \emph{The window sets the count quantum.} The earlier estimator counts edges over the window, so the speed step of one count is $\frac{60}{T_w N}$ and halves when the window doubles (Table~\ref{tab:ppr-est}). Doubling the window from 10 to 20~ms lowers the mean noise by only 20 to 34~\%, not by 50~\%, which shows that quantization is not the dominant noise above about 100 PPR.
\end{itemize}
With the earlier estimator, then, the speed resolution of the encoder limited the noise only up to about 100 PPR. Above that, the noise was set by the shaft-synchronous ripple and by timing noise, and the highest settings added outliers. These conclusions apply to the earlier estimator and are not a measurement of the slot-aligned one.

% TODO: repeat the PWM sweep with the FINAL CET estimator for several PPR settings (data kit, menus 5 and 6), and report here beside
% Table~\ref{tab:ppr-dct}: edges per revolution and lost edges (quad_err), the SD and the line at the rotation frequency by PPR, the
% bias against the index speed by PPR (accuracy by PPR), the fraction of loops with a new estimate against its prediction, and the CPU
% load. The bias by PPR shows what the resolution does to the accuracy under CET, which the DCT data of Table~\ref{tab:ppr-dct} cannot. That is the only
% test that validates the final choice, and it resolves two open questions: whether the outliers of
% 250 and 256 PPR belong to the setting or to the earlier estimator, and whether 256 PPR is better
% than 250 PPR.

\paragraph{Basis for the choice of 256 PPR.}
The setting of 256 PPR is selected for the following reasons, each of which follows from the datasheet and from the analysis of the estimator. It is decodable at the apparatus's top speed with a state width of 23.4~$\mu$s at 2500~rpm, nearly five sample intervals. It is one of the four resolutions with the tightest specified duty cycle, 49--51~\%, whereas the neighbouring 250 PPR is specified at 47--53~\%. It refreshes the slot-aligned estimate on every loop down to 23.4~rpm and reports speeds down to 2.34~rpm, which is the lowest update floor of the settings at or below 256 PPR. Its quadrature count of 1024 is also a power of two, which suits the integer arithmetic of the count alignment.

The settings from 384 to 512 PPR also pass the decoder criterion, but they present edge rates of 64 to 85~kHz at 2500~rpm against 43~kHz at 256 PPR, with a narrower margin, and their lower update floor is needed only if the apparatus is to be controlled below about 20~rpm. The choice against 125 PPR, whose update floor is 48~rpm, and in favour of 250 or 256 PPR, whose update floor is about 24~rpm, depends on the lowest controlled speed $n_{\text{low}}$ through $K \geq \frac{60}{n_{\text{low}} T_s}$. For a lowest controlled speed of 25~rpm this requires $K \geq 240$, which leaves 250 and 256 PPR among the available settings. The remaining difference between 250 and 256 PPR is small in every quantity of Tables~\ref{tab:ppr-spec} and~\ref{tab:ppr-est}. The earlier measurements give no noise-based reason to prefer 256 PPR over 125 PPR, since the noise had reached its floor by about 100 PPR (Table~\ref{tab:ppr-dct}); the case for 256 PPR rests on the update floor of the slot-aligned estimator, as stated above.

% TODO: state the lowest speed the apparatus is actually controlled at, and the lowest setpoint
% used in the N_Re-N_Po sweep. This is the number that decides between 125 PPR and 250/256 PPR.
% TODO: once the measured comparison above is complete, add a sentence stating whether it
% confirms the choice, and which of the analytical predictions it supports.

\subsection{Quadrature Decoding}
\label{subsec:agit-quad}

Both channels are registered with the pigpio daemon for callbacks on either edge, so every transition of either channel produces a callback carrying the pin, the new level, and a timestamp in microseconds. Decoding is performed by a state-transition lookup. The previous and current states of the (A, B) pair form a four-bit index into a sixteen-entry table, given in Table~\ref{tab:quad}, that returns $+1$ or $-1$ for a valid single-bit transition and $0$ for no change or for an invalid transition in which both channels appear to have changed. An invalid transition indicates that an edge was missed, and it contributes nothing to the count rather than an arbitrary $\pm 2$ step that would corrupt the position. Every such event is also counted and the count is logged with every run, in the column \texttt{quad\_err}, so that a run with a non-zero value is one in which edges were lost. The edge loss that Equation~\ref{eq:kmax} predicts for the high resolutions can thus be measured directly and not inferred. The number of callbacks delivered (\texttt{edge\_total}) is logged beside it, and over one revolution it must equal the counts per revolution.

\begin{table}[htbp]
    \centering
    \caption[Quadrature transition table]{Quadrature transition table used by the decoder (increment added to the position count). States are written as (A, B). Positive increments follow the sequence 00 $\rightarrow$ 10 $\rightarrow$ 11 $\rightarrow$ 01 $\rightarrow$ 00, in which A leads B. The diagonal entries 00$\leftrightarrow$11 and 01$\leftrightarrow$10 are invalid and contribute zero.}
    \label{tab:quad}
    \begin{tabular}{@{}c|cccc@{}}
        \toprule
        Previous $\backslash$ Current & 00 & 01 & 10 & 11 \\
        \midrule
        00 & 0 & $-1$ & $+1$ & 0 \\
        01 & $+1$ & 0 & 0 & $-1$ \\
        10 & $-1$ & 0 & 0 & $+1$ \\
        11 & 0 & $+1$ & $-1$ & 0 \\
        \bottomrule
    \end{tabular}
\end{table}

Callbacks from the pigpio daemon are delivered in order from a single thread. This property is relied upon in two places: the position count read inside the index handler is exactly the count at the instant of the index edge, and the update of the position count and of its aligned snapshot is never interleaved with another callback.

\subsection{Index Validation}
\label{subsec:agit-index}

The index channel provides one pulse per revolution, one state wide, which at 1500 rpm is 39~$\mu$s and therefore within the capability of the sampling described above. A floating or noisy index line can nevertheless produce spurious edges, and a missed pulse would silently lose a revolution, so each index edge is validated against the quadrature count accumulated since the previously accepted edge. The validation uses two thresholds derived from Equation~\ref{eq:cpr}:
\begin{itemize}
    \item An edge arriving with four counts or fewer of movement since the last accepted edge is the same index mark seen again. If the shaft has reversed across the mark this is counted as one revolution in the opposite direction; otherwise it is a double edge and is ignored.
    \item An edge arriving with fewer than 0.75 of a revolution (768 counts) of movement is rejected as noise and tallied.
    \item An edge arriving after $n > 1$ revolutions records all $n$ revolutions, since the quadrature count establishes them independently, and tallies $n - 1$ missed index pulses.
\end{itemize}
Validating against the quadrature count, rather than trusting the index line alone, makes the revolution counter robust to both noise and loss on the index channel. The internal pull-up on the index input was added after a floating line was observed to count noise as revolutions.

Each accepted index pulse is stored with its time stamp, and the log is written at the end of every run. The interval $\Delta t$ between consecutive pulses gives a speed
\begin{equation}
    n_{\mathrm{idx}} = \frac{60}{\Delta t}
    \label{eq:idx-speed}
\end{equation}
that is independent of the quadrature decoding and of the resolution setting, and the number of A/B counts between consecutive pulses must equal the counts per revolution, $4K$, which checks the DIP switch against the firmware constant.

% TODO: report the rejected and missed index counts observed over a representative run (the
% backend exposes both counters), and the count over the same run with the pull-up disabled if
% that data exists.

\subsection{Slot-Aligned Speed Estimation}
\label{subsec:agit-cet}

\paragraph{Method.}
Speed is estimated from two quantities that are both measured: the change in position count and the change in timestamp between two snapshots. This is the M/T method of \citet{bonert1983}, which combines the frequency-based and period-based methods reviewed in Section~\ref{subsec:rrl-encoder}. With the timestamp $t$ in microseconds,
\begin{equation}
    n = \frac{\Delta c}{N} \cdot \frac{60 \times 10^{6}}{\Delta t}
    \label{eq:speed}
\end{equation}
where $n$ is the speed in rpm, $\Delta c$ is the count difference, $N = 1024$ is the counts per revolution of Equation~\ref{eq:cpr}, and $\Delta t$ is the timestamp difference in microseconds. The sign of $\Delta c$ gives the direction. The timestamp subtraction is performed on unsigned 32-bit values, so that the wrap of the microsecond counter, which occurs after about 71.6 minutes, does not disturb the result.

The refinement of \citet{anuchin2019} applied here is that the count increments entering Equation~\ref{eq:speed} are restricted to multiples of four. A timestamped snapshot of the position count is stored only when the count is a multiple of four, so every pair of snapshots spans a whole number of encoder periods and always begins and ends at the same phase of the quadrature cycle. Duty-cycle and phase errors of the encoder, shown in Fig.~3 of that work, then cancel within each increment instead of appearing as speed ripple.

The control loop reads the most recent snapshot at its own scheduled instant and subtracts the previously used one. This is the ready-on-demand behaviour of the synchronous constant elapsed time method: the estimate is formed when the controller needs it, from timestamps that were captured at the edges themselves. The count of the aligned snapshot and its time stamp are stored together in one 64-bit word, so that a read by the loop can never combine a new count with an old time stamp, which would otherwise produce a single wrong estimate whose relative size is about $\frac{4}{\Delta c}$ and which is largest at low speed. Because the timestamps come from the pigpio sampler and not from the instant at which the loop runs, scheduling jitter of the loop, of the kind reported by \citet{adam2021computers}, changes which snapshot pair is used but does not corrupt the interval over which the speed is computed. The platform has no hardware quadrature peripheral or capture unit; the time stamp of every edge instead comes from the pigpio sampler, whose resolution is 5~$\mu$s rather than that of a processor clock. The comparison with the original algorithm is made in the next paragraph.

\paragraph{Synchronous operation and estimate delay.}
\citet{anuchin2019} identify the weakness of the conventional constant elapsed time (CET) method as being event driven: the estimate is produced when the encoder pulse that completes the chosen increment arrives, so it reaches the speed controller at an instant unrelated to the controller's own schedule. The estimate then has an age that varies between $T_s/2$ and $3T_s/2$, and this variable delay in the feedback forces a lower controller gain. Their remedy is to make the estimate ready on demand: the position and the time stamps of the last edges are kept, and the speed is computed immediately before the controller runs, which bounds the delay to about $T_s/2$.

The firmware has the same division of work. The pigpio callbacks, which are event driven, only record the position count and the time stamp of the edge; they compute no speed. The speed is computed in the control loop, at the loop's own instant, from the most recent aligned snapshot and the one used on the previous loop. The estimate therefore does not depend on when a pulse arrives, which removes the event-driven jitter that the method of \citet{anuchin2019} addresses. Table~\ref{tab:cet-compare} compares the three arrangements.

\begin{table}[htbp]
    \centering
    \caption[Conventional CET, synchronous CET and the MIXR-1 estimator]{Comparison of the conventional CET method, the synchronous CET method of \citet{anuchin2019}, and the estimator of this work. $T_s$ is the loop period and $T_a = \frac{60}{nK}$ is the aligned period.}
    \label{tab:cet-compare}
    \small
    \begin{tabular}{@{}p{0.15\textwidth}p{0.22\textwidth}p{0.24\textwidth}p{0.24\textwidth}@{}}
        \toprule
        & Conventional CET & Synchronous CET & MIXR-1 \\
        \midrule
        Speed computed & At the pulse that completes the increment & Before the controller runs & In the control loop, every $T_s = 10$~ms \\
        Increment $\Delta c$ & Multiple of 4, adapted to hold the elapsed time & Multiple of 4, corrected by Table~I & Multiple of 4 by construction (snapshots only when the count is a multiple of 4) \\
        Time stamp source & Timer at the event & Hardware capture units & pigpio sampler, 5~$\mu$s \\
        End of the interval & The completing pulse & The most recent edge of any type & The last aligned edge, one edge in four \\
        Delay of the estimate & $T_s/2$ to $3T_s/2$ & About $T_s/2$ & $T_s/2$ to $T_s/2 + T_a$ \\
        Increment below 4 counts & Not used & Counts 1 to 3, or a longer elapsed time & Longer elapsed time \\
        Zero speed & Timeout & Zero increment & Timeout, 100~ms \\
        \bottomrule
    \end{tabular}
\end{table}

The delay of the estimate in this firmware follows from the geometry of the snapshots. Let the loop read at time $t_r$, and let the latest aligned snapshot have been taken at $t_k \leq t_r$, so that its age $a_k = t_r - t_k$ lies between 0 and $T_a$. The interval over which the speed is averaged ends at $t_k$ and has length $\Delta t \approx T_s$, so it is centred $\Delta t/2$ before $t_k$ and the delay of the estimate relative to the read is
\begin{equation}
    d = a_k + \frac{\Delta t}{2}, \qquad \frac{T_s}{2} \leq d < \frac{T_s}{2} + T_a
    \label{eq:delay}
\end{equation}
Table~\ref{tab:delay} evaluates this for the 10~ms loop. Above about 100~rpm the delay is within 2.3~ms of the $T_s/2 = 5$~ms of the synchronous method of \citet{anuchin2019}, and it is narrower than the 5 to 15~ms range of the conventional method. The variation of the delay from one loop to the next, which is what the controller experiences as jitter, is at most $T_a$.

\begin{table}[htbp]
    \centering
    \caption[Delay of the speed estimate]{Delay of the speed estimate for $K = 256$ and $T_s = 10$~ms, from Equation~\ref{eq:delay} for the aligned estimator and Equation~\ref{eq:delay-newest} for the newest-edge estimator. The mean assumes the age is uniformly distributed over $[0, T_a)$.}
    \label{tab:delay}
    \begin{tabular}{@{}rrrrr@{}}
        \toprule
        & & \multicolumn{2}{c}{Aligned} & Newest edge \\
        \cmidrule(lr){3-4} \cmidrule(lr){5-5}
        Speed (rpm) & $T_a$ (ms) & Mean (ms) & Range (ms) & Range (ms) \\
        \midrule
        25 & 9.38 & 9.7 & 5.0--14.4 & 5.0--7.34 \\
        50 & 4.69 & 7.3 & 5.0--9.7 & 5.0--6.17 \\
        100 & 2.34 & 6.2 & 5.0--7.3 & 5.0--5.59 \\
        500 & 0.47 & 5.2 & 5.0--5.5 & 5.0--5.12 \\
        1500 & 0.16 & 5.1 & 5.0--5.2 & 5.0--5.04 \\
        \bottomrule
    \end{tabular}
\end{table}

The firmware implements two estimators that differ only in where the interval ends, and one of them is selected when the daemon starts (the option \texttt{--estimator}). The aligned estimator described above ends the interval at the last aligned edge, whose age can be up to one full period. The newest-edge estimator is the synchronous CET of \citet{anuchin2019}. It keeps the count and the time stamp of the most recent edge of each of the four types, ends the interval at the newest edge of any type, whose age is at most one state, $T_a/4$, and starts it at the most recent edge of the same type as it was at the previous estimation, so that the increment is a whole number of periods. If the increment is not a multiple of four, an edge was lost, and the estimator resynchronizes on the next call. With it the delay is
\begin{equation}
    \frac{T_s}{2} \leq d < \frac{T_s}{2} + \frac{T_a}{4}
    \label{eq:delay-newest}
\end{equation}
and the two estimators share the elapsed-time behaviour: in both, the loop period fixes the elapsed time and the increment adapts with speed, which has the same effect as the adaptation of the increment in the original method. Table~\ref{tab:delay} lists both. The cost of the newest-edge form is four time-stamped registers written on every edge instead of one on every fourth edge, which does not change the number of callbacks. Which estimator the measurements of this chapter use is decided by the comparison of the two at one resolution, with the age of the estimate logged.

% TODO: state which estimator was selected and why, from the comparison of the two at 256 PPR (age of the
% estimate, the SD of the raw speed at each step, quad_err, and the closed-loop behaviour). Everything
% measured after that comparison must use the selected estimator.
% TODO: measure the age instead of predicting it. Run with --log-age: the columns age_us and tick_rtt_us give, once
% per loop, the age of the snapshot in use (the daemon clock minus the snapshot time stamp; it includes about half
% of the round trip of the clock request, tick_rtt_us). Report its histogram at three speeds (for example 100,
% 500 and 1500 rpm) for each estimator against the distributions assumed in Table~\ref{tab:delay} (uniform on
% [0, T_a) for the aligned estimator, on [0, T_a/4) for the newest-edge estimator), and the loop-to-loop variation.

\paragraph{Window.}
Let $T_s$ denote the interval between consecutive snapshots as read by the loop. In the control loop it is the loop period, 10~ms (100~Hz). The firmware also carries a configurable RPM sample window with a default of 20~ms. That parameter belongs to an earlier version of the estimator; the slot-aligned estimator takes its interval from the aligned snapshots and does not read it. Both 10~ms and 20~ms are tabulated below so that the effect of the loop period on the error is visible.

% TODO: RPM_SAMPLE_WINDOW_US (default 20000) is still set from the command line and the network
% interface but is not read by kinematics.cpp. Remove it, or wire it in, and delete the unused
% column from Table~\ref{tab:esterr} once the final window is fixed.

\paragraph{Error analysis.}
The three classical methods have different error behaviour, which determines why the M/T form is the appropriate choice for this speed range. Following \citet{anuchin2019}, the frequency-based method counts $\Delta c$ over a fixed window and has an absolute error that is independent of speed,
\begin{equation}
    \Delta n_f = \frac{60}{T_s N}
    \label{eq:fm}
\end{equation}
the period-based method times a fixed increment of four counts with a timing resolution $h$ and has an error that grows with the square of speed,
\begin{equation}
    \Delta n_p = \frac{h \, n^{2} N}{60 \cdot 4}
    \label{eq:tm}
\end{equation}
and the M/T method, in which the count is exact and only the elapsed time is uncertain, has an error proportional to speed over the window,
\begin{equation}
    \Delta n_{MT} = \frac{h}{T_s}\, n
    \label{eq:mt}
\end{equation}
where $h = 5\,\mu$s is the effective timestamp uncertainty, the sample interval of the pigpio sampler. Table~\ref{tab:esterr} evaluates the three expressions for this encoder, with $N = 1024$.

\begin{table}[htbp]
    \centering
    \caption[Speed-estimation error of the three methods]{Absolute speed-estimation error (rpm) of the frequency-based, period-based, and M/T methods for $N = 1024$ and $h = 5\,\mu$s, and the aligned-snapshot period $T_a = \frac{60}{nK}$ at each speed. Frequency-based and M/T values are shown for $T_s = 10$ and $20$~ms.}
    \label{tab:esterr}
    \begin{tabular}{@{}rrrrrrr@{}}
        \toprule
        & & \multicolumn{2}{c}{Frequency-based} & Period-based & \multicolumn{2}{c}{M/T (this work)} \\
        \cmidrule(lr){3-4} \cmidrule(lr){6-7}
        $n$ (rpm) & $T_a$ (ms) & 10 ms & 20 ms & & 10 ms & 20 ms \\
        \midrule
        50 & 4.69 & 5.86 & 2.93 & 0.05 & 0.025 & 0.013 \\
        100 & 2.34 & 5.86 & 2.93 & 0.21 & 0.050 & 0.025 \\
        500 & 0.47 & 5.86 & 2.93 & 5.33 & 0.250 & 0.125 \\
        1000 & 0.23 & 5.86 & 2.93 & 21.33 & 0.500 & 0.250 \\
        1500 & 0.16 & 5.86 & 2.93 & 48.00 & 0.750 & 0.375 \\
        \bottomrule
    \end{tabular}
\end{table}

Two conclusions follow. At low speed the frequency-based error is unacceptable, 5.86~rpm at a 10~ms window and 2.93~rpm at 20~ms, which is 12~\% and 6~\% of a 50~rpm setpoint. At high speed the period-based error is unacceptable, 48~rpm or 3.2~\% at 1500~rpm. The M/T error is the smallest of the three across the whole range and is a constant fraction of the reading, $\frac{h}{T_s}$, which is 0.05~\% for a 10~ms window and 0.025~\% for 20~ms. The estimator therefore has a resolution that is limited by the 5~$\mu$s sampler rather than by the count quantization, which is why a high-resolution encoder setting is unnecessary. The frequency-based columns of the table describe the earlier DCT estimator, which counts edges over a window (Section~\ref{subsec:agit-cetdct}).

The M/T error accounts for timestamp uncertainty only. The residual contribution of the encoder's own position error remains. If the $0.25^{\circ}$ rated accuracy is taken as a once-per-revolution component, its effect on the speed estimated over one window is bounded by $0.25^{\circ} \times \pi / 180^{\circ} = 0.44$~\% of the reading, independent of speed and window length. This bound rests on the assumption that the error varies slowly over a window, which the datasheet does not state; the measured scatter in Section~\ref{subsec:agit-speed-test} is the test of it.

\paragraph{Delay and low-speed behaviour.}
The delay of the estimate is given by Equation~\ref{eq:delay}; at 100 rpm it is at most about 7~ms and at 1000~rpm about 5~ms, small against the mechanical time constant of the impeller.

At low speed the aligned period exceeds the loop period. A new snapshot is available on every loop only while $T_a \leq T_s$, which holds above
\begin{equation}
    n_{\text{upd}} = \frac{60}{K \, T_s} = \frac{60}{256 \times 0.010} = 23.4 \text{ rpm}
    \label{eq:nupd}
\end{equation}
Below this speed the loop sometimes sees no new aligned snapshot, and the previous estimate is held until one arrives, so the estimate is piecewise constant with steps every $T_a$. The held value is exact when it updates, since at low speed the timestamp error $\frac{h}{T_a}$ becomes negligible; what degrades is its freshness. A shaft that produces no new aligned snapshot for 100~ms is declared stopped and the speed is set to zero, which is the stopped-shaft requirement identified in Section~\ref{subsec:rrl-encoder}. This sets the minimum steady speed that can be reported as non-zero at one aligned step per 100~ms,
\begin{equation}
    n_{\min} = \frac{60}{K \times 0.100} = 2.34 \text{ rpm}
    \label{eq:nmin}
\end{equation}
When the shaft restarts after a stop the previous snapshot is not reset, so the first estimate spans the stopped interval and is close to zero; the next estimate, one aligned period later, is correct, so the restart costs one aligned period of delay and no error afterwards.

\subsection{Spectral Analysis of the Speed Signal}
\label{subsec:agit-fft}

The standard deviation of the speed in a steady hold is a single number that does not say where the scatter comes from. A scatter made of independent noise is spread evenly over frequency, whereas an error that repeats every revolution appears as lines at the rotation frequency and its multiples. This section separates the two with the Fourier transform of the logged speed.

\paragraph{Method.}
The speed is logged at 100~Hz, so the spectrum extends to 50~Hz. For each constant-speed segment the mean is removed, a Hann window is applied, and the amplitude spectrum is computed. The amplitude of the line at harmonic $h$ of the rotation frequency $f_{rot} = \frac{n}{60}$ is taken from the power in a band of $\pm 0.5$~Hz around $h f_{rot}$. Harmonics above 50~Hz are aliased to $|h f_{rot} - k \times 100\ \text{Hz}|$ and are used only when the aliased frequency is more than 4~Hz from zero, where it can be told apart from slow drift. The remainder, called broadband noise below, is the root-mean-square value of the spectrum between 1 and 48~Hz outside $\pm 1$~Hz of the first four harmonics.

Because each estimate is a mean over a window $T_s$, a component at true frequency $f$ is attenuated by $\left|\dfrac{\sin(\pi f T_s)}{\pi f T_s}\right|$, and the line amplitudes are divided by this factor wherever an angle is derived from them.

\paragraph{An angular error that repeats every revolution.}
If the encoder reports the shaft angle with an error $e(\theta) = \varepsilon_h \sin(h\theta)$, for example because the shaft runs eccentrically relative to the encoder, the speed computed from two snapshots a window $\Delta t$ apart is in error by $\dfrac{e(\theta_k) - e(\theta_{k-1})}{\Delta t}$. For a speed $n$ this is a ripple at frequency $h f_{rot}$ with relative amplitude
\begin{equation}
    \frac{A_h}{n} = h\, \varepsilon_h \left| \frac{\sin(\pi h f_{rot} T_s)}{\pi h f_{rot} T_s} \right|
    \label{eq:ripple}
\end{equation}
which is independent of the speed apart from the attenuation factor. A true ripple of the shaft speed, produced by a torque that varies once per revolution, behaves differently: it is filtered by the inertia of the shaft, and its amplitude in rpm falls as the speed rises. The two can therefore be distinguished by how the line amplitude varies with speed.

\paragraph{Observation.}
Figure~\ref{fig:orders-spectra} shows the spectrum of the raw speed in the open-loop sweep at four speeds, for the two estimators. At every speed there are narrow lines that sit on $f_{rot}$, $2f_{rot}$, and $3f_{rot}$, and Figure~\ref{fig:orders-map} shows that they follow these frequencies as the speed changes. They are present in both estimators with the same amplitude. Above them lies a broadband floor.

\begin{figure}[htbp]
    \centering
    \includegraphics[width=\textwidth]{fig_orders_spectra}
    \caption[Spectra of the raw speed at four speeds]{Amplitude spectrum of the raw speed in the open-loop PWM sweep at four speeds, for the CET and DCT estimators. The dotted lines mark the first three harmonics of the rotation frequency; above 50~Hz they are drawn at their aliased frequency.}
    \label{fig:orders-spectra}
\end{figure}

\begin{figure}[htbp]
    \centering
    \includegraphics[width=0.85\textwidth]{fig_orders_map}
    \caption[Order map of the open-loop sweep]{Amplitude spectrum of each step of the open-loop CET sweep against the speed of the step, with the first three harmonics of the rotation frequency overlaid (aliased above 50~Hz).}
    \label{fig:orders-map}
\end{figure}

\paragraph{Quantification.}
Table~\ref{tab:orders-cet} gives the standard deviation, the line amplitudes, and the share of the variance that the lines carry for the CET estimator, and Figure~\ref{fig:orders-amp} shows the amplitudes and the equivalent angular error.

\begin{table}[htbp]
    \centering
    \caption[Lines at the rotation frequency, open-loop CET sweep]{Spectral decomposition of the raw speed in the open-loop PWM sweep with the CET estimator. $A_1$, $A_2$, and $A_3$ are the amplitudes of the lines at one, two, and three times the rotation frequency, $\varepsilon_1$ is the equivalent angular error of Equation~\ref{eq:ripple} for the first line, and the share is the fraction of the variance carried by the three lines. A dash marks a line that aliases onto slow drift.}
    \label{tab:orders-cet}
    \small
    \resizebox{\textwidth}{!}{%
    \begin{tabular}{@{}rrrrrrrrr@{}}
        \toprule
        Speed (rpm) & SD (rpm) & $A_1$ (rpm) & $A_2$ (rpm) & $A_3$ (rpm) & $\frac{A_1}{n}$ (\%) & $\varepsilon_1$ (deg) & Share (\%) & Noise (rpm) \\
        \midrule
        227 & 3.76 & 1.53 & 3.61 & 0.18 & 0.67 & 0.387 & 55 & 2.38 \\
        492 & 2.02 & 1.26 & 1.61 & 0.22 & 0.26 & 0.148 & 52 & 1.17 \\
        757 & 1.51 & 1.11 & 0.75 & 0.23 & 0.15 & 0.086 & 40 & 0.53 \\
        1021 & 1.36 & 0.80 & 0.44 & -- & 0.08 & 0.047 & 23 & 1.10 \\
        1281 & 1.59 & 0.80 & 0.24 & 0.21 & 0.06 & 0.039 & 15 & 1.23 \\
        1545 & 1.15 & 0.87 & -- & 0.13 & 0.06 & 0.036 & 29 & 0.75 \\
        1805 & 1.60 & 0.98 & 0.24 & 0.13 & 0.05 & 0.036 & 20 & 0.77 \\
        2065 & 1.79 & 1.27 & 0.27 & -- & 0.06 & 0.043 & 26 & 1.21 \\
        2325 & 1.92 & 1.68 & 0.23 & 0.20 & 0.07 & 0.054 & 40 & 1.22 \\
        2499 & 2.96 & 1.10 & 0.67 & 0.15 & 0.04 & 0.034 & 10 & 2.53 \\
        \bottomrule
    \end{tabular}}
\end{table}

\begin{figure}[htbp]
    \centering
    \includegraphics[width=\textwidth]{fig_orders_amplitude}
    \caption[Line amplitude and equivalent angular error]{Left: amplitude of the lines at one and two times the rotation frequency against speed, for the CET and DCT estimators. Right: equivalent angular error of the first line (Equation~\ref{eq:ripple}) against speed, with the rated accuracy of the encoder.}
    \label{fig:orders-amp}
\end{figure}

\paragraph{Interpretation.}
Three things follow from the data.
\begin{enumerate}
    \item \emph{The scatter is largely periodic in the shaft rotation.} Between 10 and 55~\% of the variance of the raw speed, depending on the speed, sits in the lines at one to three times the rotation frequency, and the first line alone is 0.8 to 1.7~rpm at every speed. This part of the scatter is not noise of the estimator. It repeats every revolution, which is what an eccentric shaft produces, and it has the same amplitude in the CET and DCT estimators because it comes from the shaft and not from the arithmetic.
    \item \emph{It does not depend on the encoder resolution.} The line amplitude is the same at every resolution from 48 to 256 PPR (Table~\ref{tab:ppr-dct}), as an error of the shaft angle is independent of how finely the angle is counted.
    \item \emph{Two mechanisms are present.} Above about 1000~rpm the relative amplitude of the first line is flat at 0.04 to 0.08~\% (Table~\ref{tab:orders-cet}), which is the signature of an angular error independent of speed (Equation~\ref{eq:ripple}); its equivalent angle is 0.03 to 0.05$^{\circ}$, about a fifth of the $0.25^{\circ}$ rated accuracy of the encoder. Below 1000~rpm the amplitude in rpm is largest and falls with speed, and the equivalent angle rises to $0.39^{\circ}$, which is the behaviour of a true speed ripple rather than of an angular error. At 227~rpm the line at twice the rotation frequency, 3.6~rpm, is larger than the first.
\end{enumerate}
The data are therefore consistent with an eccentric shaft but do not isolate the encoder as the origin: they bound the part attributable to angular error of the encoder at about $0.05^{\circ}$ above 1000~rpm, and they show that the larger low-speed ripple behaves as a speed ripple of the shaft.

\paragraph{Ripple against shaft angle.}
The spectrum places the ripple at multiples of the rotation frequency but does not show where in the revolution it occurs. For speeds up to 800~rpm the loop sees at least seven samples per revolution, and each sample can be placed at its shaft angle with the index pulses: the angle is the time since the previous index pulse divided by the time between the previous and the next pulse. After the mean is removed, the speed is fitted by the first and second harmonics of the shaft angle. A ripple that belongs to the shaft or to the mounting is locked to the angle, so its amplitude and its phase against the index are the same revolution after revolution. The phase measured in this way also contains the delay $d$ of the speed estimate, which turns the ripple by $360^{\circ} f_{rot} d$, so the phase of the first harmonic is a straight line in the rotation frequency whose slope gives $d$ and whose intercept at zero speed gives the angle at which the shaft makes the speed highest. The slope is an independent measurement of the delay of Equation~\ref{eq:delay}, for each estimator. Two further observations separate an angular error of the encoder from a true speed ripple. An angular error gives a relative amplitude $\frac{A_1}{n}$ that does not depend on the speed (Equation~\ref{eq:ripple}), and the equivalent angle must stay below the rated $0.25^{\circ}$ of the encoder if the encoder is the cause. A true speed ripple of the shaft gives an amplitude in rpm that falls as the speed rises. Finally, the radial runout of the shaft at the encoder, measured with a dial indicator, predicts the angular error as the runout divided by the sensing radius, and can be compared with the equivalent angle found from the spectrum.

% TODO: the angle-locked amplitude and phase at each speed, the phase slope (the delay of each estimator) and the zero-speed phase (data kit,
% menu 3, 'RIPPLE VERSUS SHAFT ANGLE'), the equivalent angle against speed (menu 3, 'SHAFT / ECCENTRICITY'), and the dial-indicator runout
% in mm with the sensing radius of the AMT102-V, so that predicted and measured angle can be compared.

\paragraph{Consequences.}
The ripple sets a floor on the standard deviation that neither a finer encoder nor a better estimator can lower. At 1500~rpm the CET estimator has a standard deviation of 1.15~rpm in the open-loop sweep, of which the first line alone is 0.87~rpm. Because the ripple is periodic with the rotation, it cancels exactly in any average taken over a whole number of revolutions, and the index channel provides the revolution count needed to do so. A speed average over whole revolutions would therefore remove this ripple from the steady-state values used in the $N_{Re}$--$N_{Po}$ sweep, at no cost in accuracy. The speed controller sees the ripple as a disturbance at the rotation frequency, 25~Hz at 1500~rpm.

% TODO: separate the encoder error from the speed ripple directly. Measure the radial runout of the
% shaft at the encoder with a dial indicator, and compute the predicted angular error e/r for
% the sensing radius r of the AMT102-V; compare it with the 0.03-0.05 degrees found above.
% TODO: state whether the sweeps were run with the impeller and fluid or unloaded. A two-bladed
% impeller would put a true speed ripple at twice the rotation frequency, which is the line that
% dominates at low speed. State the number of blades.
% TODO: the 2x line aliases onto DC drift above about 1500 rpm and is not reported there. Logging
% the speed at a higher rate for a short run (for example 500 Hz) would show the full spectrum.
% TODO: the torque channel is zero in every log supplied, so no torque spectrum is available.
% When the torque sensor is logged, repeat the analysis on it: a line at the rotation frequency
% in the torque would confirm the speed ripple as a true load effect.

\subsection{Comparison of the CET and DCT Estimators}
\label{subsec:agit-cetdct}

The two estimators were compared in the daemon itself, with the encoder at 256 PPR (DIP 0011, 1024 counts per revolution), real-time (FIFO) scheduling, a 10~ms loop, and no torque channel, so that the estimator is the only thing that changes between runs. Five variants were run in turn: the DCT estimator with sampling windows of 5, 10, and 20~ms, and the two CET estimators of Section~\ref{subsec:agit-cet} (aligned and newest edge). Each variant ran three tests with the same schedules: an open-loop PWM staircase from 10 to 90~\% in steps of 10~\%, held 10~s each (test A), and closed-loop staircases of setpoints from 10 to 1500~rpm held 15~s each, once with the earlier gains $K_p = 1.5641$ and $K_i = 41.2249$ (test Bold) and once with the gains designed on the identified plant, $K_p = 0.2482$ and $K_i = 6.8691$ (test Bimc; Section~\ref{subsec:agit-pi-tuning}). The interior of every hold is analysed. The speed of every step is also computed from the index pulses (Section~\ref{subsec:agit-speed-test}) and is used as the reference for the bias. Each variant was run once, so a difference of 10 to 20~\% between two variants is inside the run-to-run scatter and is not claimed as a difference; the differences reported below as findings are factors of 3 to 11.

\paragraph{The DCT estimator.}
The earlier discrete-time (DCT) estimator was computed in the control loop. On every loop it read the position count and the system clock and formed the difference of each from the previous loop. It accumulated the count difference and the elapsed time until the elapsed time reached a configured window $T_w$, and then computed
\begin{equation}
    n = \left| \frac{\Delta c}{N} \right| \cdot \frac{60}{\Delta t}
    \label{eq:dct}
\end{equation}
with $\Delta c$ the accumulated count and $\Delta t$ the accumulated clock time. The value was held until the next window completed. This is the frequency-based method of Equation~2 in \citet{anuchin2019}, with the time taken from the loop clock. Three properties distinguish it from the CET estimator of Section~\ref{subsec:agit-cet}.
\begin{enumerate}
    \item The count increment is any integer and not a multiple of four, so each window carries one count of quantization, a speed step of $\frac{60}{T_w N}$, which is 11.7, 5.9, and 2.9~rpm for 5, 10, and 20~ms, and the duty-cycle and phase errors of the encoder are not cancelled.
    \item The elapsed time is the difference of two readings of the loop's own clock and not the time stamps of the edges. The count is read slightly before the clock, so a delay between the two readings enters the speed as a relative error equal to the delay divided by the window: 100~$\mu$s is 2~\% of a 5~ms window and 0.5~\% of a 20~ms window.
    \item The window is a threshold on the accumulated time, so the interval between refreshes is a whole number of loop periods and depends on the exact period of each loop. A window equal to a whole number of nominal periods is reached or missed depending on whether the periods are a little below or above nominal, and the refresh then waits for the next loop.
\end{enumerate}

\paragraph{Open loop.}
Table~\ref{tab:ol-sd} and Figure~\ref{fig:est-sd}a give the standard deviation of the raw speed in the open-loop staircase. The aligned CET estimator has 2.7 to 4.9~rpm from 230 to 2140~rpm, nearly constant over the range, and the newest-edge estimator is within 0.5~rpm of it at every step. The M/T error of Equation~\ref{eq:mt} is 0.1 to 0.5~rpm over this range (Table~\ref{tab:esterr}), so most of the CET scatter is the shaft-synchronous ripple of Section~\ref{subsec:agit-fft}. The DCT scatter grows with the speed, in proportion to it: from 7 to 8~rpm at 230~rpm to 51, 50, and 32~rpm at 2140~rpm for the windows of 5, 10, and 20~ms. That is 2.4, 2.3, and 1.5~\% of the speed, an equivalent timing uncertainty of 0.12, 0.23, and 0.30~ms for the window in question. This is the order of magnitude of the clock-reading delay of item~2 above and not of the count quantization of item~1, which does not depend on the speed. At 1300~rpm the DCT scatter is 31, 31, and 19~rpm against 3.6~rpm for aligned CET, a factor of 8.6, 8.6, and 5.3. Averaged over the steps from 30 to 80~\% PWM (770 to 2140~rpm) it is 34.5, 33.7, and 21.1~rpm for DCT at 5, 10, and 20~ms against 3.8~rpm for each CET estimator. The 20~ms window is the quietest DCT variant because it has the smallest count quantum and averages over the longest time, at the price of the longest delay.

\begin{table}[htbp]
    \centering
    \caption[Open-loop speed scatter by estimator]{Standard deviation of the raw speed (rpm) in the open-loop PWM staircase at 256 PPR, 10~s per step, one run per estimator. The 90~\% step is not usable for DCT: the 5 and 10~ms runs reached the 12\,000~rpm safety limit during that step (the 5~ms run before its interior could be analysed), and the 20~ms run contains a loop stall of 1.8~s.}
    \label{tab:ol-sd}
    \small
    \resizebox{\textwidth}{!}{\input{Chapter5/Figs/table_open_loop_sd}}
\end{table}

\begin{figure}[htbp]
    \centering
    \includegraphics[width=\textwidth]{fig_est_sd}
    \caption[Speed scatter against speed for five estimators]{Standard deviation of the raw speed for the three DCT windows and the two CET estimators at 256 PPR. (a) Open-loop staircase against speed. (b) Closed loop with the gains from the identified plant, against setpoint. (c) Closed loop with the earlier gains. One run per estimator.}
    \label{fig:est-sd}
\end{figure}

\paragraph{Spikes.}
The DCT estimator also produces isolated estimates far from the true speed. Counting the raw samples that differ from the local median by more than 50~\% at speeds above 100~rpm gives 851, 4524, and 1795 per million samples for DCT at 5, 10, and 20~ms, and 524 and 518 per million for the aligned and newest-edge CET estimators; the largest sample is 8.8 times the median for DCT at 10~ms and 3.6 times for CET. The same kind of sample ended two open-loop runs, DCT at 5 and 10~ms, at the 12\,000~rpm limit during the 90~\% step, where the shaft turns at about 2400~rpm. The CET figure is not zero, about one sample in two thousand, so no estimator is free of outliers, but DCT at 10~ms produces nine times as many as CET and DCT at 20~ms about three times as many. A lost edge was counted in three of the five variants (\texttt{quad\_err} equal to one for DCT at 5 and 10~ms and for newest-edge CET), so the spikes are not caused by lost counts. One plausible cause, which these logs do not test, is that the count read by the DCT estimator is updated when the callbacks are delivered and not when the edges occur, so a delayed delivery is followed by a burst that the estimator divides by a clock interval that did not stretch with it.

\paragraph{Accuracy against the index pulses.}
Scatter is not accuracy. Table~\ref{tab:est-acc} gives the mean of each estimator over each step, minus the speed from the index pulses, for speeds above 100~rpm. Every estimator is unbiased to within 0.1~\% except in the steps in which it spiked. The mean absolute bias is 0.002 to 0.008~\% for the aligned CET estimator, 0.002 to 0.018~\% for the newest-edge CET estimator, 0.007~\% for DCT at 5~ms, 0.03 to 0.06~\% for DCT at 20~ms, and 0.09 to 0.48~\% for DCT at 10~ms, where the last figure is dominated by a bias of $-3.8$~\% at the 90~\% step of test A. The DCT estimator at 10~ms is the only variant with a systematic bias, $-0.05$ to $-0.11$~\% in most steps; the cause is not identified. By Equation~\ref{eq:npo-unc} a speed bias of 0.1~\% is 0.2~\% in the Power number, small against the uncertainty of the torque (Section~\ref{subsec:agit-power}). The choice of estimator therefore changes the \emph{scatter} of the speed, and with it the behaviour of the control loop, much more than it changes the mean speed used in the $N_{Re}$--$N_{Po}$ curve.

\begin{table}[htbp]
    \centering
    \caption[Bias of each estimator against the index pulses]{Bias of the mean speed of each estimator against the speed from the index pulses, over the steps above 100~rpm. Tests: A, open loop; Bold, closed loop with the earlier gains; Bimc, closed loop with the gains from the identified plant. ``Steps'' is the number of steps analysed; the bias is in percent of the speed except the last column (rpm).}
    \label{tab:est-acc}
    \small
    \resizebox{\textwidth}{!}{\input{Chapter5/Figs/table_accuracy}}
\end{table}

\paragraph{Age of the estimate.}
The firmware logs the age of the snapshot in use, the daemon clock minus the time stamp of the snapshot, with half of the round trip of the clock request (about 50~$\mu$s) subtracted. Table~\ref{tab:age} gives it against the maximum that the geometry of Section~\ref{subsec:agit-cet} predicts, $T_a$ for the aligned and $T_a/4$ for the newest-edge estimator. The newest-edge estimator has the smaller age at every speed: by 46~\% at 100~rpm (median 1.00 against 1.85~ms), 28~\% at 250~rpm, 15~\% at 500~rpm, and 1 to 4~\% at 1000 and 1500~rpm, where $T_a$ is below 0.25~ms and no longer fixes the age. The measured age exceeds the predicted maximum for the newest-edge estimator at every speed and for the aligned estimator from 250~rpm. At 1000 and 1500~rpm the median is 0.74 to 0.78~ms for both estimators, against predicted maxima of 0.04 to 0.23~ms. The excess is the same for both estimators, so it is a delay that does not come from the estimator: the transit of the edge through the pigpio notification stream to the daemon, which the time stamps do not contain. It is small against the 5~ms of $T_s/2$ in Equation~\ref{eq:delay}. The 95th percentile of the age at 500~rpm is 1.43~ms for the aligned and 1.13~ms for the newest-edge estimator. At 350~rpm the aligned run has one age of 10.9~ms, a loop that was served late, the same kind of event as the late wake-ups of the loop timing of Section~\ref{sec:plat-timing}.

\begin{table}[htbp]
    \centering
    \caption[Measured age of the speed estimate]{Age of the snapshot in use (microseconds) for the two CET estimators, from the log of the daemon. The prediction is the maximum of the geometric argument of Section~\ref{subsec:agit-cet}, and the aligned period is $T_a = \frac{60\times10^{6}}{256\,n}$~$\mu$s. ``rtt'' is the round trip of the clock request.}
    \label{tab:age}
    \small
    \resizebox{\textwidth}{!}{\input{Chapter5/Figs/table_age}}
\end{table}

The delay was also estimated from the slope of the phase of the first harmonic of the ripple, folded against the shaft angle with the index pulses (Section~\ref{subsec:agit-fft}). The slope gives 15.7 and 14.8~ms for the two CET estimators and 29.4 and 27.0~ms for DCT at 10 and 20~ms, but $-7$~ms for DCT at 5~ms, which is not possible, and it rests on three speeds, 230 to 780~rpm, where the line is at most 2.5~rpm. It is not used here; the age of Table~\ref{tab:age} is the measurement of the delay.

\paragraph{Closed loop.}
With the PI controller closed on the estimate (Table~\ref{tab:cl-sd}, Figure~\ref{fig:est-sd}b and~c), the estimators are indistinguishable at 100~rpm and below and differ a great deal above. At setpoints from 10 to 100~rpm every estimator has a standard deviation of 8 to 27~rpm with either set of gains, except DCT at 20~ms at 50~rpm with the earlier gains (45~rpm), so that scatter belongs to the motor and its drive at low speed and not to the estimator; Section~\ref{subsec:agit-pi-tuning} treats that range. From 250~rpm the CET estimators hold the speed within 2 to 5.5~rpm at every setpoint with either set of gains, while the scatter of DCT grows with the setpoint: with the gains of the identified plant it is 8 to 33~rpm for DCT at 5~ms, 8 to 34~rpm at 10~ms, and 5 to 20~rpm at 20~ms. At 1500~rpm the DCT scatter is 32.6, 34.3, and 20.1~rpm against 3.05~rpm for aligned CET, which is 10.7, 11.2, and 6.6 times larger. Averaged over the setpoints from 250 to 1500~rpm the ratio is 5.7, 5.6, and 3.3. With the earlier, higher gains the DCT scatter at 1500~rpm is 40, 46, and 28~rpm, a ratio of 14, 16, and 9.7, so the lower gains reduce the DCT scatter by 20 to 28~\% but do not remove it, and leave the CET scatter unchanged. That the CET scatter does not depend on the gains and the DCT scatter does only in part points to the DCT estimate itself, with its spikes and its irregular refresh, and not to the tuning. These logs do not separate the contribution of the delay from that of the noise of the estimate, but the DCT window of 5~ms, which refreshes on every loop and adds the least delay of the three, is not quieter than the 20~ms window, which indicates that the noise of the estimate and not its delay dominates the scatter in this test.

\begin{table}[htbp]
    \centering
    \caption[Closed-loop speed scatter by estimator, identified-plant gains]{Standard deviation of the raw speed (rpm) in the closed-loop staircase at 256 PPR with the gains from the identified plant, $K_p = 0.2482$ and $K_i = 6.8691$, 15~s per setpoint, one run per estimator. The results for the earlier gains, $K_p = 1.5641$ and $K_i = 41.2249$, are in Figure~\ref{fig:est-sd}c.}
    \label{tab:cl-sd}
    \small
    \input{Chapter5/Figs/table_closed_loop_Bimc}
\end{table}

\paragraph{Which estimator is used.}
The aligned CET estimator is selected. From 250~rpm upward its scatter is the smallest or within a few tenths of a rpm of the smallest in every test, it has the smallest bias against the index pulses, no lost edge in any run, and a bounded delay. The newest-edge estimator is equivalent in scatter and bias and has the smaller age, but the difference in age is a few tenths of a millisecond at the speeds of the experiment (above 250~rpm) and each variant was run once, so the logs do not distinguish the two. All measurements of the following sections use the aligned estimator with the loop period of 10~ms, and the DCT estimator is not used for any measurement of the $N_{Re}$--$N_{Po}$ curve. The conclusion that CET is better than DCT at the same encoder setting, loop, and gains rests on factors of 3 to 11 in closed loop and of 5 to 9 in open loop above 1000~rpm, far larger than the run-to-run scatter, and it holds for every DCT window tried.

\paragraph{Earlier logs.}
Two earlier sets of logs are superseded by the comparison above. The 256 PPR sweeps logged by the dashboard on 2~September (DCT) and on 4 and 7~September (CET) gave a mean standard deviation of 40 to 56~rpm with DCT and 4.9 to 8.6~rpm with CET, the same order as Table~\ref{tab:ol-sd}, but they were taken on different days and under different conditions. The harness sweeps of the earlier estimator, 1 to 3~rpm for DCT, disagree with the 8 to 54~rpm of the daemon, and a harness closed-loop staircase with DCT at 20~ms refreshed the estimate on only 39~\% of the loops; the new runs, made in the same daemon with the same encoder setting, replace them.


\subsection{Output Filtering}
\label{subsec:agit-filter}

The estimate of Equation~\ref{eq:speed} is smoothed by a first-order exponential moving average before it is displayed or logged,
\begin{equation}
    \bar{n}_k = \alpha\, n_k + (1 - \alpha)\, \bar{n}_{k-1}, \qquad \alpha = 0.15
    \label{eq:ema}
\end{equation}
evaluated once per loop, $T_s = 10$~ms. The corresponding time constant, cut-off frequency, and variance reduction for uncorrelated input noise are
\begin{equation}
    \tau = -\frac{T_s}{\ln(1 - \alpha)} = 61.5\ \text{ms}, \qquad f_c = \frac{1}{2\pi\tau} = 2.6\ \text{Hz}, \qquad \frac{\sigma_{\mathrm{out}}^{2}}{\sigma_{\mathrm{in}}^{2}} = \frac{\alpha}{2 - \alpha} = 8.1\ \%
    \label{eq:ema-props}
\end{equation}
so the filter reaches 95~\% of a step in 18.4 samples, 184~ms, and the standard deviation of uncorrelated noise falls to 28~\%. The filter is re-seeded from the first reading after a reset so that a restart does not begin from a stale value.

The displayed and logged speed is the filtered value. The speed controller applies its own feedback filter, described in Section~\ref{subsec:agit-pi-structure}, so the two are separate filters with separate purposes.

% TODO: Figure 5.3 -- measured raw and filtered speed at a constant setpoint (for example 500 rpm)
% over about 10 s, annotated with the raw-speed standard deviation. Compare that standard
% deviation against the M/T prediction in Table~\ref{tab:esterr} (0.25 rpm at 500 rpm, 10 ms) to
% show how much of the observed scatter is timestamp quantization and how much is something
% else (encoder eccentricity, shaft vibration, scheduling).

\subsection{Validation Against a Handheld Tachometer}
\label{subsec:agit-speed-test}

The accuracy of the speed estimate is tested against a reference that does not depend on the estimator. Precision and accuracy are different properties: the standard deviation of the raw speed measures its scatter, and only a comparison of its mean with an independent speed says whether the mean is right. Two references are used.

The first is the index channel. At a constant speed, $N - 1$ revolutions elapse between the first and the last of $N$ index pulses inside the analysed window, and the mean speed is
\begin{equation}
    \bar{n}_{\mathrm{idx}} = \frac{(N-1)\cdot 60\times10^{6}}{t_N - t_1}
    \label{eq:idx-mean}
\end{equation}
in rpm, with the pulse times in microseconds. These times are stamped by the same pigpio sampler, but the result depends on none of the quadrature decoding, the resolution setting, the alignment of the snapshots, or the window of the DCT estimator, and an average over whole revolutions also removes the shaft-synchronous ripple of Section~\ref{subsec:agit-fft}, which cancels exactly in such an average. Its uncertainty is that of two time stamps of 5~$\mu$s over a window of several seconds, below $10^{-5}$ of the speed. The second reference is a handheld tachometer, read three times per hold, which is independent of the electronics but has a resolution and an accuracy of its own that bound the agreement the test can claim.

For each estimator and each speed the bias is the mean of the estimate over the window minus the reference, in rpm and in percent. It is reported for the CET estimators and for the DCT estimator at each window, in open loop and in closed loop, because a bias that appears only in closed loop would point to the feedback and not to the estimator. A constant bias in rpm points to an offset, a bias proportional to speed to a scale error (the counts per revolution, or a DIP switch that is not at 256 PPR), and a bias that depends on the window to the timing of the estimate. The bias translates into the Power number through Equation~\ref{eq:npo-unc}.

% TODO: Table of the bias and the scatter of each estimator at each speed, against the index reference (menu 3 of the data kit,
% 'ACCURACY'), and Figure of the bias against speed. Add the tachometer readings of the PPR campaign (test B) at 256 PPR with the
% resolution and the rated accuracy of the tachometer, and say whether the two references agree.
% TODO: the scatter against the prediction of Table~\ref{tab:esterr} (0.25 rpm at 500 rpm, 10 ms) shows how much of the observed
% scatter is timestamp quantization and how much is the shaft-synchronous ripple of Section~\ref{subsec:agit-fft}.


%********************************** 5.3 **************************************
\section{Motor Driver and Drive Logic}
\label{sec:agit-drive}

\subsection{Motor and Driver Selection}
\label{subsec:agit-motor-sel}

% TODO: state the DOGA motor ratings (voltage, rated power, rated speed, gear ratio, rated and
% stall torque) and the driver ratings (continuous and peak current, PWM frequency range).
% NOTE: reconcile the motor rating against the 80-100 W / 0.5-1.0 N*m figure quoted in the
% proposal, which does not match the as-built motor. Chapter 1 and the BOM must agree with
% whatever is stated here.

\subsection{Hardware PWM Drive Logic}
\label{subsec:agit-pwm}

Duty is generated on a hardware PWM channel rather than in software, so the waveform is unaffected by the scheduling jitter discussed in Section~\ref{sec:plat-timing}. Direction is fixed by the two direction lines, and the enable line gates the output. The initialization order that prevents a stale duty from being applied is described in Section~\ref{subsec:plat-safestart}.

The motor is driven at 20~kHz, above the audible range, on the hardware PWM channel of GPIO 13. The duty is commanded in 4095 counts, a 12-bit value that is scaled onto the 0 to $10^6$ duty range of the pigpio library. In the open-loop test modes the command is applied as given, so that a step of the duty reaches the driver within one loop period. In the interactive modes the command passes through the power ceiling and a rate limit that lets the duty rise by at most 200~\% per second, so a full-scale step takes 0.5~s; the identification of the plant in Section~\ref{subsec:agit-pi-tuning} was therefore made in a test mode, where the rate limit does not apply.

% TODO: justify the frequency against the electrical time constant of the motor (measure the current ripple
% at 20 kHz on the scope), and state the real duty resolution of the hardware channel at 20 kHz, which is
% set by the PWM clock and is not the 12 bits of the command.

\subsection{Driver Output Characterization}
\label{subsec:agit-driver-test}

Duty was swept from 0 to 100~\% in steps of 10~\% on the free motor, and through the dashboard sequences on the 96~mm and 174~mm impellers in the filled vessel (best run of each, Section~\ref{subsec:agit-plausibility}). Figure~\ref{fig:duty-rpm} gives the duty needed for each speed, the electrical and shaft power, and their ratio.

\begin{figure}[htbp]
    \centering
    \includegraphics[width=\textwidth]{fig_duty_rpm}
    \caption[Duty, power, and efficiency against speed]{(a) Duty against the steady speed for the free motor (open-loop sweep) and for the two impellers in the filled vessel. (b) Electrical power at the supply and shaft power (torque times speed). (c) Ratio of shaft to electrical power; the points at which the torque is below 0.01~N\,m are omitted.}
    \label{fig:duty-rpm}
\end{figure}

On the free motor the speed is proportional to the duty, 25.7~rpm per percent (2500~rpm at 100~\%), with an intercept of $-13$~rpm, so there is no dead band that the sweep resolves. In the filled vessel the same speed needs more duty because of the drag of the liquid. The 96~mm impeller needs 4.2~\% for 67~rpm, 35~\% for 680~rpm, and 91~\% for 1360~rpm. The 174~mm impeller needs 4.6~\% for 48~rpm, 49~\% for 280~rpm, and 75~\% for 331~rpm. The lowest duty held is therefore 4.2 to 4.6~\%, below which the sweep has no step, so the breakaway duty is below that value and is not isolated here. The steep rise of the duty with the speed for the 174~mm impeller is why its sequence ends at 350~rpm and the 96~mm sequence at 1400~rpm. The electrical power at the supply reaches 59~W for the 96~mm impeller at 1362~rpm and 70~W for the 174~mm impeller at 331~rpm, which is the quantity that the power limiter of Section~\ref{subsec:plat-powerlimit} bounds. The ratio of shaft to electrical power is 53 to 54~\% for the 96~mm impeller above 900~rpm and 16 to 21~\% for the 174~mm impeller above 150~rpm, in which the motor delivers a large torque at a low speed, so that losses in the windings probably dominate; this is not tested here. % TODO: the breakaway duty and the current ripple at 20 kHz (scope) are not measured.


%********************************** 5.4 **************************************
\section{PI Speed Control}
\label{sec:agit-pi}

\subsection{Controller Structure}
\label{subsec:agit-pi-structure}

The controller is proportional-integral, with no derivative term, since the speed measurement is differentiated from position and a derivative action would amplify its quantization noise.

Three implementation details distinguish it from the textbook form. The controller is fed the raw estimate of the speed, not the exponentially smoothed value that is displayed, and it low-pass filters that input itself with a first-order filter of coefficient $a$ (\texttt{FEEDBACK\_ALPHA}, default 0.15, a time constant of about 60~ms at the 10~ms loop), re-seeded from the live reading whenever the controller is reset so that a restart does not produce a transient from a stale filter state. The integral accumulates over the measured time step rather than an assumed period, for the reason given in Section~\ref{sec:plat-timing}. And the integral state is limited to the range that the output can cover, $|K_i I| \leq u_{\max}$, where $u_{\max}$ is the current output ceiling.

This is integrator clamping in the sense of limiting the state, and it differs from conditional integration~\citep{bohn1995}, which stops integrating while the output is saturated. It bounds the windup to one output range but does not remove it: while the output is saturated the integral can still grow up to the bound, and must unwind from it when the error changes sign. The effect on the response is measured in the power-limited run of Section~\ref{subsec:agit-pi-test}.

The ceiling $u_{\max}$ is not a constant. It is the lower of the water-safe maximum duty and the ceiling supplied by the power limiter of Section~\ref{subsec:plat-powerlimit}; in the interactive modes it is further limited to the present duty plus the allowance of the rate limit. When the limiter lowers the ceiling, the bound on the integral falls with it at the next call and the integral is pulled down at once to what the actuator is permitted to deliver, so that the controller does not hold a demand that the safety function forbids. This coupling between the safety function and the controller is the main design point of this section.

The control law, as implemented at every loop $k$ with the measured time step $\Delta t_k$, is
\begin{align}
    \tilde{n}_k &= a\, n_k + (1 - a)\, \tilde{n}_{k-1}, \qquad e_k = r_k - \tilde{n}_k \label{eq:pi-fb}\\[2pt]
    I_k &= \operatorname{clip}\!\left(I_{k-1} + e_k\, \Delta t_k,\; -\frac{u_{\max}}{K_i},\; \frac{u_{\max}}{K_i}\right) \label{eq:pi-int}\\[2pt]
    u_k &= \operatorname{clip}\!\left(K_p\, e_k + K_i\, I_k,\; 0,\; u_{\max}\right) \label{eq:pi-law}
\end{align}
where $n_k$ is the raw speed estimate of Section~\ref{subsec:agit-cet}, $a$ the feedback filter coefficient, $r_k$ the setpoint, $u_k$ the duty in counts of the 12-bit channel, and $u_{\max}$ the ceiling described above. In the continuous-time form that is used for the design below, $u(s) = C(s)\,e(s)$ with
\begin{equation}
    C(s) = K_p + \frac{K_i}{s} = K_p\left(1 + \frac{1}{T_i\, s}\right), \qquad T_i = \frac{K_p}{K_i}
    \label{eq:pi-tf}
\end{equation}
The proportional gain $K_p$ (counts per rpm) is the immediate reaction to the present error: an error of 100~rpm moves the duty by $100\,K_p$ counts at once. The integral gain $K_i$ (counts per rpm per second) accumulates the error, so the duty keeps changing until the error is zero, which is what removes the steady-state error against the drag of the fluid. The ratio $T_i = K_p/K_i$ is the time that the integral action takes to repeat the proportional action, and its inverse is the frequency at which the controller changes from mainly integral to mainly proportional.

\begin{figure}[htbp]
    \centering
    \begin{tikzpicture}[>=stealth, node distance=14mm, every node/.style={font=\small}, blk/.style={draw, minimum width=17mm, minimum height=9mm, align=center}]
        \node (r) {$r$};
        \node[circle, draw, right=10mm of r, inner sep=1pt] (sum) {$+$};
        \node[blk, right=10mm of sum] (pi) {PI\\$K_p,\,K_i$};
        \node[blk, right=12mm of pi] (sat) {clip\\$[0,\,u_{\max}]$};
        \node[blk, right=12mm of sat] (pl) {plant\\$G(s)$};
        \node[right=10mm of pl] (y) {$n$};
        \node[blk, below=9mm of pi] (est) {estimator\\+ filter $a$};
        \node[blk, above=9mm of sat, minimum width=34mm] (lim) {power limiter,\\water-safe duty};
        \draw[->] (r) -- (sum); \draw[->] (sum) -- node[above]{$e$} (pi); \draw[->] (pi) -- (sat); \draw[->] (sat) -- node[above]{$u$} (pl); \draw[->] (pl) -- (y);
        \draw[->] (lim) -- (sat); \draw[->] (lim.west) -| node[pos=0.25, above]{$u_{\max}$} (pi.north);
        \draw[->] ([xshift=5mm]pl.east) |- (est.east); \draw[->] (est.west) -| node[pos=0.25, below]{$\tilde n$} node[pos=0.95, left]{$-$} (sum.south);
    \end{tikzpicture}
    \caption[Block diagram of the speed controller]{Block diagram of the speed loop. The ceiling $u_{\max}$ comes from the power limiter and the water-safe duty and acts twice: it clips the output and it bounds the integral state of Equation~\ref{eq:pi-int}, so the integral is pulled down at once when the limiter lowers the ceiling.}
    \label{fig:pi-block}
\end{figure}

\subsection{Tuning}
\label{subsec:agit-pi-tuning}

\paragraph{Plant identification.}
The plant from the PWM command, in counts of the 12-bit channel, to the measured speed was identified from the open-loop PWM sweeps. In each sweep the command is raised in steps of 10~\% and the speed response of each step is fitted by a first-order response with dead time,
\begin{equation}
    n(t) = n_0 + K \Delta u \left[ 1 - e^{-A (t - \theta)} \right], \qquad t > \theta
    \label{eq:plant}
\end{equation}
where $\Delta u$ is the step in counts, $K$ the steady-state gain in rpm per count, $A$ the pole in s$^{-1}$, and $\theta$ the dead time, with $t = 0$ at the first row of the step, which already carries the new command. Table~\ref{tab:plant} gives the result for the CET sweep. The gain is constant at 0.64~rpm per count from 819 to 3685 counts, which is the speed gain of the motor and the driver, and the time constant is 30 to 33~ms and the dead time 8~ms at every step. The DCT sweep gives the same values within 5~\% (time constant 30.1~ms, dead time 8.3~ms), so the dead time does not depend on the estimator. It is about the delay of a CET estimate, which is half the loop period plus the age of the snapshot, 5 to 6~ms (Section~\ref{subsec:agit-cet}), plus a few milliseconds of driver latency. The sample-and-hold of the controller, which holds each command for one 10~ms loop period, is not in it, and adds half a period, 5~ms, to the delay of the closed loop. The step to 4095 counts is limited by the speed and power caps (gain 0.43~rpm per count) and is excluded, and the first step from rest is excluded because it includes the stiction of the shaft.

\begin{table}[htbp]
    \centering
    \caption[Identified plant, CET open-loop sweep]{Plant identified from the open-loop PWM sweep with the CET estimator, by Equation~\ref{eq:plant}. $K$ is the steady-state gain, $\tau = 1/A$ the time constant, $\theta$ the dead time, and RMS the residual of the fit.}
    \label{tab:plant}
    \small
    \begin{tabular}{@{}lrrrr@{}}
        \toprule
        PWM step (counts) & $K$ (rpm/count) & $\tau$ (ms) & $\theta$ (ms) & RMS (rpm) \\
        \midrule
        409 -- 819 & 0.646 & 31.3 & 8.7 & 2.0 \\
        819 -- 1\,228 & 0.645 & 32.5 & 7.7 & 1.8 \\
        1\,228 -- 1\,638 & 0.640 & 30.3 & 8.5 & 1.4 \\
        1\,638 -- 2\,047 & 0.639 & 30.8 & 8.1 & 1.6 \\
        2\,047 -- 2\,457 & 0.640 & 30.9 & 8.1 & 1.2 \\
        2\,457 -- 2\,866 & 0.636 & 30.6 & 8.4 & 1.6 \\
        2\,866 -- 3\,276 & 0.636 & 31.1 & 8.3 & 1.8 \\
        3\,276 -- 3\,685 & 0.631 & 38.6 & 5.9 & 2.3 \\
        \bottomrule
    \end{tabular}
\end{table}

The plant was identified a second time on the new open-loop sweeps, logged by the daemon with the CET estimators. With the aligned estimator the fit gives $K = 0.667$~rpm/count, $\tau = 32.9$~ms, and a dead time of 8.0~ms, and with the newest-edge estimator $K = 0.667$~rpm/count, $\tau = 35.4$~ms, and 7.7~ms. This is the same plant as in Table~\ref{tab:plant} within 10~\%, and the rule of Equation~\ref{eq:imc} gives, for $\omega_c = 5$~rad/s, $K_p = 0.237$ and $K_i = 7.21$ from the first fit and $K_p = 0.256$ and $K_i = 7.22$ from the second, against 0.233 and 7.51 from the earlier sweep. The differences are inside the uncertainty of the identification.

\begin{figure}[htbp]
    \centering
    \includegraphics[width=\textwidth]{fig_plant_id}
    \caption[Identification of the plant from the open-loop sweep]{Identification of the plant from the open-loop PWM sweep logged with the aligned CET estimator, with the time origin at the first row of each step. (a) One step (819 to 1229 counts) with the fit of Equation~\ref{eq:plant}; the dead time $\theta$ and the time $\theta + \tau$ are marked. (b) Steady-state gain $K$ and (c) time constant $\tau = 1/A$ and dead time $\theta$ of every step against the speed before the step.}
    \label{fig:plant-id}
\end{figure}

The refit of the steps from 410 to 2867 counts of this sweep, by Equation~\ref{eq:plant}, gives a gain of 0.66 to 0.67~rpm per count, a time constant of 32 to 37~ms, and a dead time of 7.5 to 8.7~ms for every step (Figure~\ref{fig:plant-id}), so the plant is linear in the sense of the design: the gain does not depend on the speed over this range and the dynamics do not change from step to step.

The 100~Hz log limits the accuracy of this identification. The time constant is only three to four samples, and the result depends on the treatment of the dead time: a first-order fit without it returns a pole of 20 to 25~s$^{-1}$, depending on the time origin, against 32~s$^{-1}$ with it, and a residual of about 3~rpm against 1.2 to 2.0~rpm. The pole and the dead time should therefore be read to about $\pm 20$~\%. The result also depends on the time origin: taking $t = 0$ at the last sample before the step, instead of at the first row of the step, adds one loop period, 10~ms, to the dead time. An earlier version of the analysis did so and found 18~ms; the comparison with measured closed-loop responses below exposed the error.

\paragraph{Derivation of the gain rule.}
Take the identified plant of Equation~\ref{eq:plant}, whose transfer function from the duty in counts to the speed in rpm is
\begin{equation}
    G(s) = \frac{K A}{s + A}\, e^{-\theta s}
    \label{eq:plant-tf}
\end{equation}
and the controller of Equation~\ref{eq:pi-tf}. The controller has one zero, at $s = -\dfrac{K_i}{K_p}$, and the plant has one pole, at $s = -A$. The design places the zero on the pole,
\begin{equation}
    \frac{K_i}{K_p} = A
    \label{eq:cancel}
\end{equation}
so that the slow part of the plant, its time constant $1/A \approx 31$~ms, is cancelled by the controller and what remains of the open loop $L(s) = C(s)\,G(s)$ is an integrator with a delay,
\begin{equation}
    L(s) = K_p\,\frac{s + A}{s}\cdot\frac{K A}{s + A}\, e^{-\theta s} = \frac{K_p K A}{s}\, e^{-\theta s}
    \label{eq:loop}
\end{equation}
The gain crossover $\omega_c$ is the frequency at which $|L(j\omega_c)| = 1$. For Equation~\ref{eq:loop} it is
\begin{equation}
    \omega_c = K_p K A
    \label{eq:wc}
\end{equation}
and it is the one number that sets how fast the loop is. Without the delay the closed loop is
\begin{equation}
    \frac{L(s)}{1 + L(s)} = \frac{\omega_c}{s + \omega_c}
    \label{eq:cl-first}
\end{equation}
a first-order response with time constant $\dfrac{1}{\omega_c}$, whose 10 to 90~\% rise time is $\dfrac{\ln 9}{\omega_c} \approx \dfrac{2.2}{\omega_c}$ and whose 5~\% settling time is $\dfrac{3}{\omega_c}$. A crossover of 5~rad/s therefore means a rise time of 0.44~s and a settling time of 0.6~s, which the measured step response of Table~\ref{tab:cl-model} confirms (0.46~s and 0.68~s). Solving Equations~\ref{eq:cancel} and~\ref{eq:wc} for the gains gives the rule
\begin{equation}
    K_p = \frac{\omega_c}{K A}, \qquad K_i = K_p A = \frac{\omega_c}{K}
    \label{eq:rule}
\end{equation}
in counts per rpm for $K_p$ and counts per rpm per second for $K_i$. The crossover is chosen from the delay of the loop, as explained below; for a plant with dead time the internal-model-control choice $\omega_c = 1/(T_c + \theta)$, with $T_c$ the desired closed-loop time constant, gives
\begin{equation}
    K_p = \frac{1}{K A\,(T_c + \theta)}, \qquad K_i = \frac{1}{K\,(T_c + \theta)}
    \label{eq:imc}
\end{equation}
which is Equation~\ref{eq:rule} with $\omega_c = 1/(T_c + \theta)$. The designed bandwidth is $\omega_c \approx 1/T_c$ only when $T_c \gg \theta$. For $\omega_c = 5$~rad/s the rule gives $K_p = 0.233$ and $K_i = 7.51$ for the plant of Table~\ref{tab:plant}, and a check of $K_p K A$ against the target is printed with the gains.

\paragraph{What the phase margin is and why it matters.}
Equation~\ref{eq:loop} has a phase of $-90^{\circ}$ from the integrator and a further lag from the delay, which grows with frequency: a delay $\theta_{\mathrm{tot}}$ turns a component of frequency $\omega$ by $\omega\,\theta_{\mathrm{tot}}$ radians. The phase margin is the distance of the phase of the open loop from $-180^{\circ}$ at the crossover,
\begin{equation}
    \mathrm{PM} = 180^{\circ} + \angle L(j\omega_c) = 90^{\circ} - \frac{180^{\circ}}{\pi}\,\omega_c\,\theta_{\mathrm{tot}}
    \label{eq:pm}
\end{equation}
where $\theta_{\mathrm{tot}}$ is the total delay of the loop: the dead time of the plant, 8.2~ms, plus half the loop period for the sample-and-hold of the controller, 5~ms, plus any further delay of the speed estimate. The gain margin is the factor by which the loop gain can rise before the loop becomes unstable,
\begin{equation}
    \mathrm{GM} = \frac{1}{|L(j\omega_{180})|}, \qquad \angle L(j\omega_{180}) = -180^{\circ}
    \label{eq:gm}
\end{equation}
Both are measures of the same thing, how far the loop is from oscillating on its own, and the phase margin has four consequences for this apparatus.
\begin{enumerate}
    \item \emph{Stability.} At $\mathrm{PM} = 0$ the loop sustains an oscillation. The phase margin is also the extra delay that the loop can take at the crossover before that happens, the delay margin $\mathrm{PM}\,(\pi/180^{\circ})/\omega_c$.
    \item \emph{Damping.} The margin sets the overshoot and the ringing of the step response. A second-order loop with a phase margin of $60^{\circ}$ has a damping ratio near 0.6 and an overshoot of about 10~\%, and above about $70^{\circ}$ the response has almost no overshoot. As a rule of thumb the damping ratio is $\zeta \approx \mathrm{PM}/100^{\circ}$ for margins below $65^{\circ}$. The torque and the Power number are measured in steady state, so an oscillating or ringing speed would add the scatter that the averaging of Section~\ref{subsec:agit-sweep} has to remove.
    \item \emph{Robustness to the load.} The plant is not the free motor once the impeller is in the water: its gain $K$ changes by 30~\% over the speed range, and between impellers by more (Section~\ref{subsec:agit-pi-tuning}). A change of $K$ changes the loop gain in proportion, so the gain margin is the factor of change that is tolerated.
    \item \emph{Robustness to the estimator.} The estimate of the speed is itself a delay. Adding the 19~ms of a DCT estimate costs $\omega_c \times 19$~ms of phase, which is a small loss at a low crossover and a large one at a high crossover.
\end{enumerate}
The loops of this work are designed for a phase margin of at least $60^{\circ}$ and a gain margin of at least 8~dB, which are the usual limits for a loop that must stay well damped under a change of load.

Table~\ref{tab:margins2} and Figure~\ref{fig:bode-pm} give the margins of the two sets of gains with the identified free-motor plant, $K = 0.640$~rpm/count, $1/A = 31.0$~ms, and a total delay of 13.2~ms for a CET estimate and 32.2~ms for a DCT estimate. With the earlier gains the crossover is at 29.2~rad/s and the phase margin falls from $74^{\circ}$ with a CET estimate to $42^{\circ}$ with a DCT estimate, and the gain margin from 11.7~dB (a factor 3.8) to 4.5~dB (a factor 1.7). The present gains cross at 4.8~rad/s and keep $86^{\circ}$ and $81^{\circ}$, with gain margins of 27.9 and 20.1~dB. The magnitude curves of the CET and DCT cases coincide because a delay changes only the phase.

\begin{figure}[htbp]
    \centering
    \includegraphics[width=0.62\textwidth]{fig_bode_pm}
    \caption[Open-loop Bode plot and phase margin]{Open-loop frequency response $L(j\omega)$ of the speed loop (Equation~\ref{eq:loop} with the identified plant of the free motor) for the earlier and the present gains, with a CET and a DCT estimate. The crossover (magnitude 0~dB, dots) and the phase margin, the distance of the phase from $-180^{\circ}$ at the crossover, are marked. The magnitude is the same for a CET and a DCT estimate, since a delay does not change it.}
    \label{fig:bode-pm}
\end{figure}

\begin{table}[htbp]
    \centering
    \caption[Gain, phase and delay margins of the two sets of gains]{Crossover, phase margin, gain margin and delay margin of the speed loop for the earlier and the present gains, computed from the identified free-motor plant ($K = 0.640$~rpm/count, $1/A = 31.0$~ms, dead time 8.2~ms, 5~ms of sample-and-hold). The DCT estimate adds 19~ms. Model values, not measurements.}
    \label{tab:margins2}
    \small
    \begin{tabular}{@{}llrrrrr@{}}
        \toprule
        Gains & Estimate & $\omega_c$ (rad/s) & PM ($^{\circ}$) & GM (dB) & GM (factor) & Delay margin (ms) \\
        \midrule
        Earlier, $K_p = 1.564$, $K_i = 41.2$ & CET & 29.2 & 74 & 11.7 & 3.8 & 44 \\
        & DCT & 29.2 & 42 & 4.5 & 1.7 & 25 \\
        Present, $K_p = 0.233$, $K_i = 7.51$ & CET & 4.8 & 86 & 27.9 & 24.8 & 314 \\
        & DCT & 4.8 & 81 & 20.1 & 10.1 & 294 \\
        \bottomrule
    \end{tabular}
\end{table}

\paragraph{Choice of the crossover.}
Equation~\ref{eq:pm} sets a limit on the crossover for a given delay, and the speed it buys is the other side of the trade. Figure~\ref{fig:design-tradeoff} shows both. With a CET estimate the margin falls below $60^{\circ}$ only above a crossover of about 40~rad/s, and with a DCT estimate above about 16~rad/s, so the estimator sets how fast the loop can be made. Within the limit, a faster loop is quicker (a response reaches 90~\% in $\ln 10/\omega_c = 2.3/\omega_c$ seconds) and rejects a disturbance sooner. A slower loop is calmer: the shaft-synchronous ripple of Section~\ref{subsec:agit-fft} is at 25~Hz (157~rad/s) at 1500~rpm and at 4~Hz (25~rad/s) at 250~rpm, and a controller with a high crossover passes it to the duty. The crossover of 5~rad/s is a compromise for a rig that holds a speed for tens of seconds and not a servo: a rise time of 0.46~s against holds of 15 to 30~s, a margin of $86^{\circ}$ that leaves room for the change of the plant with load and speed, and a crossover well below the ripple at the lowest controlled speed.

\begin{figure}[htbp]
    \centering
    \includegraphics[width=\textwidth]{fig_design_tradeoff}
    \caption[Phase margin, step response and overshoot against the designed crossover]{Effect of the designed crossover $\omega_c$ with the gains of Equation~\ref{eq:rule}. (a) Phase margin against $\omega_c$ for a CET estimate (total delay 13.2~ms) and a DCT estimate (32.2~ms), with the $60^{\circ}$ limit. (b) Simulated response to a step of 100~rpm for five values of $\omega_c$, with a CET estimate. (c) Time to reach 90~\% of the step and the overshoot against $\omega_c$. Simulations of the identified free-motor plant.}
    \label{fig:design-tradeoff}
\end{figure}

\paragraph{Why the gains of the earlier runs differ from these.}
The gains of the earlier closed-loop runs, $K_p = 1.5641$ and $K_i = 41.2249$, were computed for a nominal bandwidth of 1~rad/s with
\begin{equation}
    K_p = \frac{\omega_c}{K}, \qquad K_i = \frac{A\,\omega_c}{K}
    \label{eq:earlier-rule}
\end{equation}
Compared with Equation~\ref{eq:rule}, the proportional gain lacks the division by $A$, so both gains are larger by the factor $A$, which is 26 to 32~s$^{-1}$ for this plant. By Equation~\ref{eq:wc} the crossover that results is not $\omega_c$ but $K_p K A = A\,\omega_c \approx 29$~rad/s, 4.6~Hz, and not 1~rad/s. For 1~rad/s the rule gives $K_p = 0.048$ and $K_i = 1.55$. The two pairs are therefore not comparable parameter by parameter: the 41 of the earlier runs is an integral gain, to be compared with 7.5 now, and its proportional gain of 1.56 with 0.23 now. The earlier pair is a high-bandwidth tuning, and the present pair is a deliberate 5~rad/s tuning. The ratio $K_i/K_p$ is the plant pole in both, 26 in the earlier pair and 32 in the present one, which is the cancellation of Equation~\ref{eq:cancel}.

With the earlier gains the loop has a phase margin of $74^{\circ}$ with a CET estimate, which is well damped, as the measured step response confirms (Table~\ref{tab:cl-model}), but only $42^{\circ}$ when the estimate is the DCT estimate refreshed every 20 to 30~ms, and $34^{\circ}$ at the longer interval, which is poorly damped. With the present gains the margin is $80^{\circ}$ to $86^{\circ}$ for every estimator, so the loop no longer depends on which estimator supplies the speed. The feedback filter of the controller with its default coefficient 0.15 would reduce the margin of the earlier gains to $33^{\circ}$ even with a CET estimate, which is evidence that the earlier runs used no filter (coefficient 1.0). Table~\ref{tab:margins} gives these cases.

\begin{table}[htbp]
    \centering
    \caption[Gain crossover and phase margin of the two sets of gains]{Gain crossover and phase margin of the loop for the earlier gains and the present gains, computed from the identified plant ($K = 0.640$~rpm/count, $\tau = 31.0$~ms, dead time 8.2~ms) with 5~ms added for the sample-and-hold of the controller. The extra delay of the DCT estimate over the CET estimate is taken as 19~ms at the mean refresh interval of 25~ms and 24~ms at the longer interval of 30~ms (an estimate delay of about 25 and 30~ms against about 6~ms for CET, from half the window plus half the refresh interval). The last column adds the feedback filter of the controller with coefficient 0.15 (time constant 61.5~ms) to a CET estimate. These are model values, not measurements.}
    \label{tab:margins}
    \small
    \begin{tabular}{@{}lrrrrr@{}}
        \toprule
        & Crossover & \multicolumn{4}{c}{Phase margin (degrees)} \\
        \cmidrule(lr){3-6}
        Gains & (rad/s) & CET & DCT mean & DCT worst & CET + filter \\
        \midrule
        Earlier: $K_p = 1.5641$, $K_i = 41.2249$ & 29.2 & 74 & 42 & 34 & 33 \\
        Present: $K_p = 0.2329$, $K_i = 7.509$ & 4.8 & 86 & 81 & 80 & 71 \\
        \bottomrule
    \end{tabular}
\end{table}

\paragraph{Check against measured closed-loop step responses.}
The identified plant was checked against two closed-loop staircases of the present firmware, logged on the free motor with the CET aligned estimator, with the integral kept at each setpoint step. One was run with the earlier gains and one with the gains that the data kit computed from its own identification. The gains of the second were not logged, and were recovered from the response of the duty to the error, which is the sum of $K_p\, e_k$ and $K_i T_s \sum e$: the same method returns the earlier gains as 1.5640 and 41.216, against 1.5641 and 41.2249, and returns $K_p = 0.2484$ and $K_i = 6.870$ for the second run. Table~\ref{tab:cl-model} gives the rise time, the overshoot, and the 5~\% settling time, averaged over the five steps of 250~rpm between 500 and 1500~rpm, measured and from a simulation of the identified plant with the controller sampled every 10~ms. The simulation of a plant with $K = 0.667$~rpm/count, a pole of 30~s$^{-1}$, and a dead time of 8~ms reproduces both responses to within 1~\% of the step (root-mean-square error 0.008 of the step for each), and the steady-state error is below 0.1~rpm at every step. With the earlier value of the dead time, 18~ms, the same simulation predicted an overshoot of 16~\% for the earlier gains where 1.6~\% was measured, and that disagreement exposed the error of one loop period in the time origin of the identification. The gains of Equation~\ref{eq:imc} are therefore verified on the free motor with the CET estimator: the loop behaves as designed, with a rise time of 0.46~s against a designed crossover of 5~rad/s, no overshoot to speak of, and no steady-state error.

\begin{table}[htbp]
    \centering
    \caption[Measured and simulated closed-loop step response]{Closed-loop step response of the free motor with the CET aligned estimator, measured and simulated, averaged over five steps of 250~rpm from 500 to 1500~rpm. The log has a resolution of 10~ms, which is the resolution of the rise times. Single repeat.}
    \label{tab:cl-model}
    \small
    \resizebox{\textwidth}{!}{%
    \begin{tabular}{@{}lrrrrrr@{}}
        \toprule
        & \multicolumn{3}{c}{Measured} & \multicolumn{3}{c}{Simulated} \\
        \cmidrule(lr){2-4} \cmidrule(lr){5-7}
        Gains & Rise (ms) & Overshoot (\%) & Settling (s) & Rise (ms) & Overshoot (\%) & Settling (s) \\
        \midrule
        $K_p = 1.5640$, $K_i = 41.22$ & 30 & 1.6 & 0.09 & 40 & 0 & 0.09 \\
        $K_p = 0.2484$, $K_i = 6.870$ & 460 & 0.9 & 0.68 & 470 & 0 & 0.63 \\
        \bottomrule
    \end{tabular}}
\end{table}

\begin{figure}[htbp]
    \centering
    \includegraphics[width=\textwidth]{fig_pi_step}
    \caption[Measured and simulated closed-loop staircases]{Closed-loop staircase of the free motor with the aligned CET estimate and the two sets of gains: measured speed (grey), setpoint (black), and a simulation of the identified plant ($K = 0.667$~rpm/count, $A = 30$~s$^{-1}$, dead time 8~ms) with the controller of Equation~\ref{eq:pi-law} at 10~ms (colour). Top: the whole staircase from 500 to 1500~rpm. Bottom: one step of 250~rpm. The earlier gains reach the setpoint in about 30~ms, the identified-plant gains in about 0.5~s.}
    \label{fig:pi-step}
\end{figure}

The earlier harness runs differed in one respect: they restarted the controller at each setpoint step, so that the duty fell to $K_p$ times the error and the speed dipped before it recovered, whereas the present staircase keeps the integral.

% TODO: repeat the check with each impeller fitted and the vessel filled (data kit, menu 4, choice 3), with the CET newest estimator, and with
% the DCT estimator; the loaded plant has a different gain and time constant, and the common pair of gains must be checked on both.
% TODO: the rise time of the earlier gains is below the 10 ms resolution of the log in part; state it as 'about 30 to 40 ms'.

\paragraph{One pair of gains for two impellers.}
The firmware holds a single pair $(K_p, K_i)$, so the plants of the two impellers, each identified with the vessel filled as in use and with the baffles in and out, must be served by the same pair. The plant also changes with the speed, because the drag of the fluid grows with it: for the 96~mm impeller in the vessel, the identification gives a gain $K$ of 0.48~rpm per count in the neighbourhood of 290~rpm and 0.34 at 800~rpm, a pole of $39.5$~s$^{-1}$ (time constant 25~ms), and a dead time of 7.7~ms, against $K = 0.667$, a pole of 30.4~s$^{-1}$, and 8.0~ms for the free motor. The product $K A$, which sets the loop gain of the controller, falls by 22~\% from the free motor to the loaded impeller, which is what a larger inertia gives, and the gain $K$ alone falls with the speed. The heavier impeller has a smaller steady-state gain and a larger time constant, so a pair designed for one of them by Equation~\ref{eq:imc} is too weak or too aggressive for the other. The pair is therefore found by search. For every candidate pair the loop $L(s) = (K_p + K_i/s)\, G_k(s)$ is evaluated for each plant $G_k$ of Equation~\ref{eq:plant}, and the pair is chosen that maximizes the gain crossover of the slower of the two loops, subject to a phase margin of at least $60^{\circ}$ and a gain margin of at least 8~dB for both, and to a crossover below 10~rad/s for both, which is well inside the 100~Hz sampling of the loop. The result is reported together with the best pair for each impeller alone, so that the cost of the compromise is stated as the fraction of the crossover that the slower loop keeps.

% TODO: state the baffle condition of the 96 mm identification above (plant_96.csv). Identify the 174 mm impeller and the other baffle
% conditions; from the steady states of the dashboard sweeps its gain K is only about 0.07 to 0.18 rpm per count, so the common pair will
% make it slower than the 96 mm impeller: report by how much.
% TODO: the plants and the common pair: Kp, Ki, the crossover, the phase margin and the gain margin of each impeller, and the closed-loop
% check of each impeller with the common pair (data kit, menu 4, choices 1 to 3). If the compromise leaves one impeller with less than
% about half of its best crossover, say so and state whether a pair per impeller is warranted, which needs the firmware to switch the pair
% with the impeller menu.
\paragraph{Effect of $K_p$, $K_i$, and the feedback filter on the loaded rig.}
The rule of Equation~\ref{eq:rule} is a design for the free motor. On the loaded rig, the 174~mm impeller with baffles in the filled vessel, the gains were examined by a matrix of 18 settings: $K_p = 2.629$ and $3.155$, $K_i = 1.454$, 6, and 12, and a feedback filter coefficient $a = 1$, 0.2, and 0.15, each started from rest to 300~rpm with a duty limit of 35~\%, then stepped to 340~rpm at 10~s and back at 14~s. Figure~\ref{fig:gain-matrix} gives the response of three of them and the metrics of every setting at $K_p = 2.629$ (one run per setting, so differences of a few tenths of a percent are not significant).
\begin{itemize}
    \item \emph{$K_i$ sets the speed of the response.} Raising it from 1.454 to 12 shortens the rise from 10 to 90~\% from 5.2--5.3~s to 1.2--1.3~s, the settling time to $\pm 5$~\% from 6.3--6.5~s to 3.9--4.1~s, and the standard deviation of the speed at 300~rpm from 2.9--3.6 to 0.5--0.7~rpm. The price is overshoot, from about 1~\% to 10--12~\% at start-up; the duty limit holds the output at its ceiling during the rise, and the integral keeps accumulating, so part of that overshoot is the windup that the clamp of Equation~\ref{eq:pi-int} limits but does not remove.
    \item \emph{The feedback filter removes the jitter of the duty and not the speed response.} With $a = 1$ the duty changes by 0.12 to 0.44~\% from one loop to the next at steady state, because it follows the quantization of the speed estimate. With $a = 0.2$ or 0.15 and $K_i \geq 6$ it does not change at the one-count resolution of the log, and at $K_i = 1.454$ it changes by 0.16 to 0.24~\%, with the same rise time, settling time, and overshoot as without the filter. The filter costs phase (Table~\ref{tab:margins}) but is worth it where the actuator has to stay quiet.
    \item \emph{$K_p$ has little effect in the range tested.} Changing $K_p$ from 2.629 to 3.155 changes the rise and settling times by less than 0.2~s and the overshoot by about one percentage point.
\end{itemize}
The gains that the firmware carries by default, $K_p = 2.629$, $K_i = 12$, and $a = 0.15$, are among the fastest settings of this matrix, with a speed scatter of 0.6 to 0.7~rpm and a quiet duty, and the overshoot of 11.6~\% is the cost. Their ratio $K_i/K_p = 4.6$~s$^{-1}$ is not the plant pole of Equation~\ref{eq:cancel}, so this is a tuning of the loaded rig by measurement and not by the rule. % TODO: state the impeller of the disturbance runs (the duty of 15 % for 350 rpm is that of the 96 mm impeller) and confirm with the run notes that these are the gains of config.hpp, and say whether the overshoot of 11.6 % is acceptable for the N_Re-N_Po sequence.

\begin{figure}[htbp]
    \centering
    \includegraphics[width=\textwidth]{fig_gain_matrix}
    \caption[Effect of the gains on the loaded rig]{Effect of $K_i$ and the feedback filter coefficient $\alpha$ on the response of the loaded rig (174~mm impeller, baffled, duty limit 35~\%) to a start from rest to 300~rpm, at $K_p = 2.629$. (a) Speed of three settings. (b) Their duty. (c) to (f): overshoot, settling time to $\pm 5$~\%, jitter of the duty (standard deviation of the change of duty per loop) and standard deviation of the speed at 300~rpm, against $K_i$ (rows) and $\alpha$ (columns); lighter is better. One run per setting.}
    \label{fig:gain-matrix}
\end{figure}
% TODO: the gains must be found again for the working load (impeller fitted and fluid in the vessel): the plant of Table~\ref{tab:plant}
% is the unloaded motor. Menu 4 of the data kit does this and reports the plant of the loaded rig.
% TODO: confirm by a closed-loop run, and report here, that the phase margins of Table~\ref{tab:margins} agree with the damping measured
% (settling time and overshoot) for the present gains.
% TODO: state the FEEDBACK_ALPHA used in cet_1500rpm.csv and dct_1500rpm.csv (1.0 expected from Table~\ref{tab:margins}).

\subsection{Closed-Loop Speed Tracking}
\label{subsec:agit-pi-test}

Two tests compare the closed loop with the open loop on the loaded rig. The speed was held at 350~rpm (the closed-loop runs) or at a constant duty that gives about 350~rpm (the open-loop run) in the filled vessel, and the load was disturbed by five drops of water onto the impeller and by inserting and removing the baffles. The closed-loop runs used two crossovers, 0.5 and 1.0~rad/s. Figure~\ref{fig:disturbance} gives the three runs, and Table~\ref{tab:disturbance} the size and the duration of the disturbances.

\begin{figure}[htbp]
    \centering
    \includegraphics[width=\textwidth]{fig_disturbance}
    \caption[Open loop and PI control against load disturbances]{Speed in the filled vessel for a constant duty (open loop, top) and for the PI controller with crossovers of 0.5 and 1.0~rad/s (middle, bottom). Dotted lines: water drops. Dashed lines: baffles inserted (orange) and removed (green). Right: the five water drops of each run overlaid, as a percentage of the speed before the drop.}
    \label{fig:disturbance}
\end{figure}

\begin{table}[htbp]
    \centering
    \caption[Disturbance rejection, open loop against PI]{Response to the load disturbances of Figure~\ref{fig:disturbance}. Dip: the largest fall of the speed within 12~s of a drop (2nd percentile of the filtered speed) or 15~s of the insertion of the baffles, in percent of the speed before. Noise: the same measure over the 3~s before the drop, which is the dip that scatter alone produces. Recovery: the time after which the 0.5~s mean of the speed stays within 5~\% of its value before. Final: the speed 12 to 20~s after the disturbance relative to before. Mean over the five drops.}
    \label{tab:disturbance}
    \small
    \begin{tabular}{@{}llrrrrr@{}}
        \toprule
        Disturbance & Control & Before (rpm) & Dip (\%) & Noise (\%) & Recovery (s) & Final (\%) \\
        \midrule
        Water drop & Open loop & 316 & 21.8 & 0.9 & not recovered & $-18.2$ \\
        & PI, 0.5 rad/s & 350 & 3.6 & 2.0 & $<0.1$ & $+0.1$ \\
        & PI, 1.0 rad/s & 350 & 2.9 & 2.1 & 0.3 & $+0.0$ \\
        Baffles in & Open loop & 225 & 34.9 & -- & not recovered & $-32.2$ \\
        & PI, 0.5 rad/s & 350 & 7.7 & -- & 13 & $+0.4$ \\
        & PI, 1.0 rad/s & 350 & 5.8 & -- & 2.6 & $-0.6$ \\
        \bottomrule
    \end{tabular}
\end{table}

In open loop a drop of water lowers the speed by 22~\% on average, and the speed does not come back: it settles 18~\% lower, since the duty is constant and the drag of the liquid is not. Inserting the baffles takes 35~\% of the speed and the speed stays 32~\% lower. Under the PI controller the dip of a drop is 3~\%, which is no larger than the 2~\% that the scatter alone produces before the drop, and the speed returns to 350~rpm in less than a third of a second; the baffles take 6 to 8~\% and the speed is back within 2.6~s at 1.0~rad/s. At 0.5~rad/s the recovery from the baffles takes 13~s, because the speed oscillates for several seconds after the disturbance (the burst in the middle panel of Figure~\ref{fig:disturbance}), which the faster loop damps sooner. The steady-state error is zero within the noise of the measurement in every closed-loop case and 18 to 32~\% in the open-loop case, which is the reason for a speed controller on this apparatus: the Power number depends on the square of the speed, so an error of 18~\% in speed would be an error of 36~\% in $N_{Po}$ (Equation~\ref{eq:npo-unc}).

The staircase of the free motor (Figure~\ref{fig:pi-step}) and the matrix of the loaded rig (Figure~\ref{fig:gain-matrix}) complete the tracking tests.

% TODO: compare the steady-state error with the open-loop knob positioning of the CEK-A (Section~\ref{subsec:che3214l}); this is the headline comparison in Chapter~\ref{ch:validation}.


%********************************** 5.5 **************************************
\section{Torque Sensing and Power Computation}
\label{sec:agit-torque}

\subsection{Sensor Selection}
\label{subsec:agit-torque-sel}

The transducer is a reaction-torque sensor rated at 2~N\,m with a flange to square-drive interface, supplied with its own transmitter. Section~\ref{sec:rrl-agitation} established why reaction torque, rather than electrical power at the motor terminals, is the appropriate basis for the Power number: the electrical measurement includes drive-train losses that are not delivered to the fluid.

% TODO: Figure 5.7 -- torque sensor installed in the shaft line, showing the coupling arrangement
% and the reaction restraint. State the rated accuracy from the datasheet and confirm that the
% expected torque range from the impeller sizing falls within the 2 N*m rating with margin.

\subsection{Digital Acquisition over Modbus RTU}
\label{subsec:agit-torque-acq}

The transmitter exposes the measurement as holding registers over RS485, read with Modbus function code~03. No external analog-to-digital converter is used; the resolution of the chain is that of the transmitter.

% TODO: Table 5.4 -- serial and register configuration as used: port, baud, framing, slave
% address, register address and count, word order, integer format, and the scale factor with its
% justification from the transmitter's decimal-point setting.

The reading path is defensive by design. The port is opened once at start-up and probed; if the probe fails, the channel reports itself unavailable and the rest of the system continues without torque rather than aborting. Each response is length- and CRC-checked before decoding, and a read that times out yields no value rather than a stale or zero one.

\subsection{Zero-Point Stability and Automatic Tare}
\label{subsec:agit-tare}

Span and coarse zero are established on the transmitter itself against known loads, which, as Section~\ref{subsec:rrl-torque-methods} notes, no software correction replaces. Residual drift is removed by an automatic software tare that re-zeroes the channel only when the rig is demonstrably idle.

Four conditions must hold simultaneously before a tare is accepted: the drive must be off, the shaft must have been stationary for a settling interval, a full window of readings must be quiet and free of trend, and the resulting zero must lie close to the transmitter's own zero. The first two are separate requirements because a shaft can be stationary with torque still applied, by a stalled motor or by fluid still swirling after a stop; the last two guard against taring out a real load.

% TODO: Table 5.5 -- the tare acceptance thresholds as configured, with the measured idle noise
% that each was derived from.

\subsection{Zero-Offset and Drift Characterization}
\label{subsec:agit-torque-drift}

% TODO: TC-U3. Log the unloaded reading from cold start over an extended period. Report Figure 5.8
% (reading vs time) and quantify warm-up drift and steady-state noise. This is the evidence base
% for the thresholds in Table 5.5 and connects to the drift behaviour discussed in
% Section~\ref{subsec:rrl-torque-methods}.

\subsection{Validation Against Theoretical Torque}
\label{subsec:agit-torque-test}

The reaction-torque sensor was checked against known torques. A mass $m$ hung at a lever arm $r$ from the axis of the sensor produces the torque
\begin{equation}
    \tau_{\mathrm{ref}} = m\, g\, r
    \label{eq:tau-ref}
\end{equation}
with $g = 9.80665$~m\,s$^{-2}$, $m$ taken from the nominal mass (100 to 1000.5~g, in steps of 100~g) and $r = 24.77$~mm, so the reference runs from 0.024 to 0.243~N\,m. The reference rests on two assumptions, that the masses are as labelled and that the lever arm is as measured; its uncertainty is therefore that of the mass and of the arm, which is not stated for the masses used, and a relative error of 0.5~\% in the arm is an error of 0.5~\% in every reference value. % TODO: state the uncertainty of the masses and of the lever arm (the calliper and the balance used), and the rated accuracy of the sensor from the datasheet.
Each load was applied five times, and the reading is the mean of 10 samples after the zero (baseline) of the trial is subtracted. Figure~\ref{fig:loadcell} and Table~\ref{tab:loadcell} give the result for the 50 trials. A straight-line fit of the measured on the applied torque has a slope of 0.9978 and an intercept of $-0.15$~mN\,m, with $R^{2} = 0.99995$. The root-mean-square error is 0.68~mN\,m and the largest error 1.3~mN\,m. The mean error is negative at every load, between $-0.05$ and $-1.2$~\%, with the largest relative errors at the two lowest loads (0.024 and 0.049~N\,m), where an absolute error of 0.3~mN\,m is 1.2~\% of the load.

\begin{figure}[htbp]
    \centering
    \includegraphics[width=\textwidth]{fig_loadcell}
    \caption[Torque sensor against known torques]{Torque sensor against the reference torque of Equation~\ref{eq:tau-ref} (50 trials, 10 loads). (a) Measured against applied torque, with the 1:1 line and the fit. (b) Residual of every trial. (c) Relative error of the mean of each load, with the standard deviation of the five trials as the bar.}
    \label{fig:loadcell}
\end{figure}

\begin{table}[htbp]
    \centering
    \caption[Torque sensor against known torques]{Mean reading of the torque sensor for each load against the reference of Equation~\ref{eq:tau-ref}, five trials per load. SD is the standard deviation of the five trial means.}
    \label{tab:loadcell}
    \small
    \begin{tabular}{@{}rrrrr@{}}
        \toprule
        Mass (g) & Reference (N\,m) & Mean (N\,m) & SD (mN\,m) & Error (\%) \\
        \midrule
        100.0 & 0.0243 & 0.0240 & 0.0 & $-1.20$ \\
        200.0 & 0.0486 & 0.0480 & 0.0 & $-1.20$ \\
        300.0 & 0.0729 & 0.0728 & 0.3 & $-0.10$ \\
        400.0 & 0.0972 & 0.0970 & 0.0 & $-0.17$ \\
        500.0 & 0.1215 & 0.1214 & 0.5 & $-0.05$ \\
        600.0 & 0.1457 & 0.1450 & 0.0 & $-0.51$ \\
        700.0 & 0.1700 & 0.1698 & 0.4 & $-0.14$ \\
        800.0 & 0.1943 & 0.1932 & 0.4 & $-0.58$ \\
        900.0 & 0.2186 & 0.2180 & 0.0 & $-0.28$ \\
        1000.5 & 0.2430 & 0.2425 & 1.0 & $-0.24$ \\
        \bottomrule
    \end{tabular}
\end{table}

The test covers 0.024 to 0.243~N\,m, 1 to 12~\% of the 2~N\,m range of the sensor, which is the range of the torque of the $N_{Re}$--$N_{Po}$ experiment (up to 0.22~N\,m for the 96~mm and 0.32~N\,m for the 174~mm impeller) but not the whole range of the sensor. The linearity and the span (slope 0.998) are good to 0.3~\% over the range tested, and above 0.07~N\,m the sensor reads 0.05 to 0.6~\% low.

\subsection{Power Computation}
\label{subsec:agit-power}

The shaft power is the product of the measured torque $\tau$ and the angular velocity $\omega$ of the shaft,
\begin{equation}
    P = \tau\,\omega = 2\pi\, n_s\, \tau, \qquad n_s = \frac{n}{60}
    \label{eq:shaftpower}
\end{equation}
with $n$ in rpm and $n_s$ in revolutions per second. The Reynolds and Power numbers of the impeller of diameter $D$ in a fluid of density $\rho$ and viscosity $\mu$ are then
\begin{equation}
    N_{Re} = \frac{\rho\, n_s D^{2}}{\mu}, \qquad
    N_{Po} = \frac{P}{\rho\, n_s^{3} D^{5}} = \frac{2\pi\,\tau}{\rho\, n_s^{2} D^{5}}
    \label{eq:npo-torque}
\end{equation}
Because the power is measured as torque times speed, the speed enters the Power number as $n_s^{-2}$ and not as $n_s^{-3}$: one factor of the speed in the power cancels one of the three in the denominator. The relative standard uncertainty of the Power number follows from the uncertainties of the torque and of the speed, taken as independent,
\begin{equation}
    \left(\frac{u(N_{Po})}{N_{Po}}\right)^{2} = \left(\frac{u(\tau)}{\tau}\right)^{2} + \left(2\,\frac{u(n)}{n}\right)^{2}
    \label{eq:npo-unc}
\end{equation}
so that a speed error of $x$~\% becomes an error of about $2x$~\% in $N_{Po}$ and of $x$~\% in $N_{Re}$. This is why the accuracy of the speed estimate of Section~\ref{subsec:agit-speed-test} matters for the experiment and not only for the controller. The density and the viscosity of water are taken at the measured fluid temperature from the correlations of Tanaka and of Vogel, valid from 0 to 40~$^{\circ}$C. The uncertainty of the speed is the bias against the index reference combined with the scatter of the mean over the averaging window, and that of the torque is the accuracy of the sensor combined with the scatter of its mean. The propagation of Equation~\ref{eq:npo-unc} is carried into the error analysis of Section~\ref{sec:val-error}.

% TODO: Table of the torque and speed contributions to u(N_Po) at each setpoint, and which of the two dominates (data kit, menu 8,
% 'DOMINANT UNCERTAINTY'); the torque accuracy from the datasheet is an input and must be stated.
% TODO: cite the correlations (Tanaka et al., 1975, for the density of water between 0 and 40 deg C; the Vogel equation for the viscosity)
% and add them to references.bib, or replace them by the property table of the textbook used in CHE 3214L.
% TODO: unit conversion and the numerical example of one setpoint, in Appendix~\ref{app:calculations}.


%********************************** 5.6 **************************************
\section{Mode 1: Standalone Local Interface}
\label{sec:agit-mode1}

\subsection{Interface Design}
\label{subsec:agit-hmi}

% TODO: Figure 5.10 -- front panel layout. State the display type and bus address, the rotary
% encoder used for input, and the function of each button.

\subsection{Menu Flow and Input Handling}
\label{subsec:agit-menu}

% TODO: Figure 5.11 -- menu state diagram. Describe detent decoding and debouncing for the input
% knob, the deadband applied to reject jitter at rest, and the display refresh strategy that keeps
% the LCD from blocking the control loop.

\subsection{Internal Logic Validation}
\label{subsec:agit-mode1-test}

% TODO: exercise every menu transition and confirm the commanded state matches the display.
% Report as a pass/fail table.


%********************************** 5.7 **************************************
\section{Mode 2: PC Dashboard and Data Logging}
\label{sec:agit-mode2}

\subsection{Dashboard Design}
\label{subsec:agit-dashboard}

% TODO: Figure 5.12 -- dashboard screenshot with the live plots and controls annotated. State the
% GUI framework, the update rate, and which telemetry field groups it enables.

\subsection{Automated Speed Sweep and Curve Generation}
\label{subsec:agit-sweep}

The sweep automates what was previously the most laborious part of the experiment. The dashboard steps the setpoint through a sequence, waits at each point for the speed and torque to settle, records the settled values, and computes the Reynolds and Power numbers from them using the fluid properties and the inherited geometry of Section~\ref{sec:arch-basis}.

To show what the choice of speed estimator does to the curve, the same sweep was also run from the daemon, without the dashboard, once for each estimator and in random order, with the torque logged beside the speed. Each setpoint is held for 30~s and the interior of the hold (from 40 to 97~\% of it) is averaged. The Reynolds and Power numbers are computed from Equations~\ref{eq:npo-torque} twice: with the speed of the estimator, and with the speed of the index pulses of Section~\ref{subsec:agit-speed-test}. The difference between the two $N_{Po}$ values is the error that the estimator puts into the curve.

% TODO: state the settling criterion and dwell time of the dashboard sweep, and how many samples are averaged per point.
% Figure 5.13 -- a generated N_Re-N_Po curve (dashboard), and the same curve from each estimator (data kit, menu 8, fig_npo.pdf).
% TODO: report the effect of each estimator on N_Po in percent (data kit, menu 8, 'EFFECT OF THE SPEED ESTIMATOR'). A speed bias of
% x percent should appear as about 2x percent, with the opposite sign, by Equation~\ref{eq:npo-unc}.

\subsection{Data Export and Logging Integrity}
\label{subsec:agit-logging}

% TODO: TC-S4 and TC-S6. State the export formats and the column schema. Verify integrity by
% comparing the record count received against the count written over a long run, and by
% re-importing an exported file and confirming the values round-trip. Report the dropped-record
% rate, if any, and its cause.


%********************************** 5.8 **************************************
\section{Agitation Validation}
\label{sec:agit-validation}

\subsection{Endurance and Uptime}
\label{subsec:agit-endurance}

% TODO: TC-S2 and TC-S3. Run continuously for the specified duration at a representative setpoint.
% Report uptime, any faults latched, the loop-period distribution over the run (ties back to
% Section~\ref{sec:plat-timing}), thermal behaviour, and torque zero drift across the run.

\subsection{Hydrodynamic Plausibility of the Generated Curves}
\label{subsec:agit-plausibility}

The curves generated by the apparatus are compared with three references: a CFD simulation of each impeller, the Power-number curves measured by students on an Armfield stirred-vessel unit for three other impellers (a stock paddle, a custom paddle, and custom blades), and the literature curve tabulated by Brown et al.\ (curve~10 of the reference, here the ``literature'' curve), which is the curve of the 96~mm impeller configuration and exists for no other. % TODO: add the BibTeX key of Brown et al. and of the Armfield student report, and give the Armfield unit and vessel geometry.
All Reynolds and Power numbers use Equations~\ref{eq:npo-torque} with $\rho = 996.3$~kg\,m$^{-3}$, $\mu = 0.854$~mPa\,s and the diameters of the dashboard, $D = 0.0665$~m for the 96~mm impeller and $0.1207$~m for the 174~mm impeller. The CFD is available as the torque at each simulated speed (55 speeds for each impeller, 0.01 to 11\,600~rpm for the 96~mm and 0.004 to 3500~rpm for the 174~mm impeller) and is converted with the same equations. The tank was baffled in every experimental run and in the simulation. The liquid level was 200~mm, with the bottom of the impeller 50~mm above the tank floor for the 96~mm impeller and 90~mm for the 174~mm impeller.

\paragraph{Data and method.}
Each impeller was run several times through the dashboard sequence (28 setpoints for the 96~mm impeller, 23 for the 174~mm impeller, 15~s per setpoint). One run per impeller is used for the comparison, the best one, chosen by a rule fixed before the comparison: the run must be complete; the torque must not fall while the speed rises at $N_{Re} \geq 5\times10^{4}$ (a collapse of the torque, as caused by power limiting or by air drawn into the liquid); and among the runs that pass, the one with the smallest coefficient of variation of $N_{Po}$ over the plateau $N_{Re} \geq 5\times10^{4}$ is taken. The best run is that of 6 November 19:17 for the 96~mm impeller (variation 1.4~\%) and that of 6 November 21:17 for the 174~mm impeller (3.7~\%). Figure~\ref{fig:runs} shows it against the other complete runs. One run, of 6 November 19:31, has a flatter plateau (3.5~\%) but is excluded because its torque falls at the highest speeds. The Power number at the plateau of the other complete runs ranges from 1.69 to 2.09 for the 96~mm and from 2.12 to 2.53 for the 174~mm impeller, so the best run is at the upper end of the 174~mm runs, and the comparison below is for that run and not for the typical run. % TODO: confirm the cause of the runs that rise slowly (96 mm: 5 Nov 22:48, 6 Nov 19:09 and 22:09) and of the torque collapse (174 mm, 6 Nov 19:45 and 19:53) from the run notes, and the torque zero offset of each run.
Torque below 0.002~N\,m, about two counts of the sensor, is treated as noise and omitted. The CFD torque is interpolated linearly in the logarithms of speed and torque to each experimental speed. The error measure is the root-mean-square error (RMSE)
\begin{equation}
    \mathrm{RMSE} = \sqrt{\frac{1}{n}\sum_{i=1}^{n}\left(y_{\mathrm{ref},i} - y_{\mathrm{data},i}\right)^{2}}, \qquad \mathrm{NRMSE} = \frac{\mathrm{RMSE}}{\overline{y}_{\mathrm{data}}}
    \label{eq:rmse}
\end{equation}
of the torque (N\,m) or of $N_{Po}$, with $y_{\mathrm{ref}}$ the reference (CFD or literature) interpolated to the Reynolds number of each point of the data, over the points in which both exist. Because the CFD is a trusted reference only in the turbulent range, the main comparison uses $N_{Re} \geq 4\times10^{4}$; the comparison over every point above the sensor floor is also given. The error bars in the figures are the standard deviation of the torque within each averaging window of the run.

\begin{figure}[htbp]
    \centering
    \includegraphics[width=\textwidth]{fig_runs}
    \caption[All complete runs and the best run]{Power number against Reynolds number for all complete runs of the 96~mm (left) and 174~mm (right) impeller. The best run (coloured) is chosen by the rule in the text; the others are grey.}
    \label{fig:runs}
\end{figure}

\paragraph{Torque and Power number against the CFD.}
Figures~\ref{fig:cfd-96} and~\ref{fig:cfd-174} compare the best run of each impeller with the CFD, point by point, and Table~\ref{tab:rmse-cfd-mixr} gives the error. In the range $N_{Re} \geq 4\times10^{4}$ the CFD underpredicts the torque of the 96~mm impeller by 27~\% on average ($\tau_{\mathrm{CFD}}/\tau_{\mathrm{exp}} = 0.73$, RMSE 0.038~N\,m) and that of the 174~mm impeller by 18~\% (0.82, RMSE 0.048~N\,m). The Power number of the experiment is 2.02 for the 96~mm and 2.56 for the 174~mm impeller, against 1.46 and 2.08 for the CFD, an RMSE of 0.59 (NRMSE 29~\%) and 0.50 (20~\%). Below $N_{Re} = 4\times10^{4}$ the experimental points scatter widely and fall below the CFD, as the torque there is a few times the resolution of the sensor (panel~d of both figures). The ratio is nearly constant above $4\times10^{4}$ for both impellers, so the CFD and the experiment have the same dependence on speed, and the disagreement is in the level of the torque, $-18$ to $-27$~\%. The CFD is baffled like the experiment, so a difference in the baffling does not explain the gap. The conditions of the run, the zero of the torque sensor, and the geometry of the simulated impeller, which are not tested here, remain as causes.

\begin{table}[htbp]
    \centering
    \caption[CFD against the best run: torque and Power number]{Error of the CFD against the best run of each impeller in torque and in Power number. Range: ``Re $\geq 4\times10^{4}$'' (Reynolds number) or ``all above sensor floor''. $n$ is the number of setpoints, $\tau_{\mathrm{CFD}}/\tau_{\mathrm{exp}}$ the mean ratio, and the overlined quantities the means over the points. RMSE and NRMSE are defined in Equation~\ref{eq:rmse}.}
    \label{tab:rmse-cfd-mixr}
    \small
    \resizebox{\textwidth}{!}{\input{Chapter5/Figs/best_cfd_vs_mixr}}
\end{table}

\begin{figure}[htbp]
    \centering
    \includegraphics[width=\textwidth]{fig_cfd_96mm}
    \caption[96 mm impeller against the CFD]{96~mm impeller, best run (filled markers, error bars: standard deviation of the torque within each averaging window) against the CFD (open circles, one per simulated speed). (a) Power number against Reynolds number over the whole range of the CFD (1 to $10^{6}$) and of the literature curve (squares, dashed), so that every point of the three is shown; the shaded band is the range of panel~b. (b) The turbulent range on a linear axis. (c) Torque against speed. (d) Ratio of the CFD torque to the measured torque; open markers: $N_{Re} < 4\times10^{4}$; the shaded band is $\pm10$~\%.}
    \label{fig:cfd-96}
\end{figure}

\begin{figure}[htbp]
    \centering
    \includegraphics[width=\textwidth]{fig_cfd_174mm}
    \caption[174 mm impeller against the CFD]{As Figure~\ref{fig:cfd-96} for the 174~mm impeller, with the CFD and the measurement only: there is no literature curve for this impeller.}
    \label{fig:cfd-174}
\end{figure}

\paragraph{CFD against the Armfield student data.}
The students' curves are for different impellers, so the comparison measures how far their curves lie from the CFD curves of this work, and not an error of the CFD. Figure~\ref{fig:armfield} shows them and Table~\ref{tab:rmse-cfd-arm} gives the RMSE over the points with $N_{Re} \geq 4000$, the lowest Reynolds number at which the CFD is used (the first point of the stock and custom paddle, at 3912, is therefore omitted). The stock paddle has $N_{Po}$ of 7.8 to 12.1 over $N_{Re} = 1.5\times10^{4}$ to $3.4\times10^{4}$, about six times the CFD, with an RMSE of 8.2 against the 96~mm and 7.7 against the 174~mm CFD. The custom paddle falls from 1.6 at $1.8\times10^{4}$ to 0.4 at $8\times10^{4}$ and the custom blades from 8.0 at $7.9\times10^{3}$ to 0.2 at $8\times10^{4}$. Their RMSE is 0.84 and 1.42 for the paddle, and 2.2 and 2.4 for the blades, against the 96 and 174~mm CFD, which is 116 to 195~\% of their mean $N_{Po}$. The RMSE of the blades is dominated by their lowest-$N_{Re}$ points: it falls to 1.1 against the 96~mm CFD when only the points above $2\times10^{4}$ are used. The students' $N_{Po}$ keeps falling with $N_{Re}$ up to $8\times10^{4}$, whereas both CFD curves and the literature curve are flat above about $10^{4}$. A Power number that falls by more than a decade over a decade of Reynolds number is not the behaviour of fully turbulent flow, and the more likely reason is a torque reading that is small against the zero offset of the measurement at the low speeds of that unit, the same cause that the MIXR-1 data show at low speed (Section~\ref{subsec:agit-torque-drift}). That explanation is not tested here.

\begin{figure}[htbp]
    \centering
    \includegraphics[width=\textwidth]{fig_armfield}
    \caption[Armfield student data against the CFD]{Armfield student data (stock paddle, custom paddle, custom blades), the CFD of the 96~mm (left) and 174~mm (right) impeller, and the literature curve, on logarithmic axes.}
    \label{fig:armfield}
\end{figure}

\begin{table}[htbp]
    \centering
    \caption[CFD against the Armfield student data]{Error of the CFD (the 96 and 174~mm impellers of this work) against the three Armfield student curves, for the points with $N_{Re} \geq 4000$. $n$ is the number of points, the overlined quantities are means over the points, and the last column is the RMSE of $\log_{10}N_{Po}$, a measure that does not weight the large low-$N_{Re}$ values.}
    \label{tab:rmse-cfd-arm}
    \small
    \resizebox{\textwidth}{!}{\input{Chapter5/Figs/rmse_cfd_vs_armfield}}
\end{table}

\paragraph{Comparison with the literature curve.}
The curve of Brown et al.\ (curve~10) is that of the 96~mm impeller configuration, so it is compared with the 96~mm impeller only; the 174~mm impeller is compared with its CFD alone (Figure~\ref{fig:cfd-174}). Panel~(a) of Figure~\ref{fig:cfd-96} shows the three on one set of axes. The literature curve falls from 12 at $N_{Re} = 3$ to a flat level of 1.5 to 1.8 above $N_{Re} \approx 100$ and stays there to $10^{6}$. The CFD follows it at the high end, with 1.8 at $N_{Re} = 10^{4}$ and 1.5 at $10^{5}$ against 1.7 and 1.8, but it rises far above it below $5\times10^{3}$: 2.9 at 5000, 11.5 at 2000, and 36 at 1000, against 1.7, 1.6, and 1.6. The curve of the measurement lies at $N_{Po} \approx 2$ over $N_{Re} = 4\times10^{4}$ to $1.2\times10^{5}$, on the same level as the other two. Table~\ref{tab:rmse-lit} gives the error over the Reynolds numbers where each series exists. The best 96~mm run is 14~\% above the literature curve (RMSE 0.31) and the CFD is 18~\% below it (RMSE 0.32), so the measurement and the CFD lie on opposite sides of the literature curve by about the same amount, and the measurement is closer to the literature than the CFD is. The Armfield curves are far from the literature curve: the stock paddle by a factor of 5.4 (RMSE 8.0), the custom paddle 59~\% below it (RMSE 1.1) and the custom blades 19~\% below it on average (RMSE 2.4; the mean bias understates the error because the large positive errors at low $N_{Re}$ cancel the negative errors at high $N_{Re}$). The 96~mm data and the CFD therefore agree with the literature to within 14 to 18~\%, and the student curves do not agree with any of them. The geometry of the 96~mm impeller is not exactly that of the configuration of the literature curve, so the offset is not attributed to the measurement. % TODO: state the geometry of the reference configuration (impeller type, D/T, C/T, number of baffles) and in which ratios the 96 mm impeller departs from it.

\begin{table}[htbp]
    \centering
    \caption[The 96 mm data, the CFD and the Armfield data against the literature curve]{Error of the best 96~mm run, the 96~mm CFD, and the Armfield curves against the literature curve (Brown et al., curve~10), interpolated in the logarithms to the Reynolds number of each point. 96~mm run and CFD: $N_{Re}\geq4\times10^{4}$ and, for the CFD, up to the highest experimental value; Armfield: $N_{Re}\geq4000$. Bias is the mean of $(N_{Po} - N_{Po,\mathrm{lit}})/N_{Po,\mathrm{lit}}$.}
    \label{tab:rmse-lit}
    \small
    \resizebox{\textwidth}{!}{\input{Chapter5/Figs/best_vs_literature_96}}
\end{table}

\paragraph{Limits of the comparison.}
Four limits apply. The comparison uses one run per impeller, the best by the rule above, and the other runs differ from it by up to 20~\% in $N_{Po}$ at the plateau; the conditions of the run (the zero of the torque sensor and the state of the liquid) are not recorded. The liquid level is the same for the two impellers but the clearance is not (50~mm and 90~mm). The Armfield and literature data were transcribed from a table, so their Reynolds and Power numbers carry the rounding of that table. And the Reynolds number range of the apparatus ends at $1.2\times10^{5}$, so the comparison with the literature curve says nothing about the laminar and transitional range, where the stock paddle and the CFD show the largest changes.


%********************************** 5.9 **************************************
\section{Results and Discussion}
\label{sec:agit-discussion}

% TODO: synthesize. Does the speed loop meet the regulation target across the range? Does the
% torque chain resolve the range the experiment requires? Where does the dominant uncertainty in
% N_Po come from, speed or torque? What behaviour is attributable to the platform's timing rather
% than to the sensors? Carry the quantitative answers into Chapter~\ref{ch:validation}.
